\documentclass[a4paper,12pt]{paper}
%DIF LATEXDIFF DIFFERENCE FILE
%DIF DEL lecture8.tex             Fri Nov 14 00:45:20 2025
%DIF ADD diffs/lecture8_new.tex   Fri Sep 25 13:57:12 2026

\usepackage{mycommands}
\usepackage{amsmath} % Ensure amsmath is loaded for align environment

\title{Lecture 19: Waveform tomography}
\author{David Al-Attar \\ Michaelmas Term, \the\year}
%DIF PREAMBLE EXTENSION ADDED BY LATEXDIFF
%DIF UNDERLINE PREAMBLE %DIF PREAMBLE
\RequirePackage[normalem]{ulem} %DIF PREAMBLE
\RequirePackage{color}\definecolor{RED}{rgb}{1,0,0}\definecolor{BLUE}{rgb}{0,0,1} %DIF PREAMBLE
\providecommand{\DIFadd}[1]{{\protect\color{blue}\uwave{#1}}} %DIF PREAMBLE
\providecommand{\DIFdel}[1]{{\protect\color{red}\sout{#1}}}                      %DIF PREAMBLE
%DIF SAFE PREAMBLE %DIF PREAMBLE
\providecommand{\DIFaddbegin}{} %DIF PREAMBLE
\providecommand{\DIFaddend}{} %DIF PREAMBLE
\providecommand{\DIFdelbegin}{} %DIF PREAMBLE
\providecommand{\DIFdelend}{} %DIF PREAMBLE
\providecommand{\DIFmodbegin}{} %DIF PREAMBLE
\providecommand{\DIFmodend}{} %DIF PREAMBLE
%DIF FLOATSAFE PREAMBLE %DIF PREAMBLE
\providecommand{\DIFaddFL}[1]{\DIFadd{#1}} %DIF PREAMBLE
\providecommand{\DIFdelFL}[1]{\DIFdel{#1}} %DIF PREAMBLE
\providecommand{\DIFaddbeginFL}{} %DIF PREAMBLE
\providecommand{\DIFaddendFL}{} %DIF PREAMBLE
\providecommand{\DIFdelbeginFL}{} %DIF PREAMBLE
\providecommand{\DIFdelendFL}{} %DIF PREAMBLE
%DIF COLORLISTINGS PREAMBLE %DIF PREAMBLE
\RequirePackage{listings} %DIF PREAMBLE
\RequirePackage{color} %DIF PREAMBLE
\lstdefinelanguage{DIFcode}{ %DIF PREAMBLE
%DIF DIFCODE_UNDERLINE %DIF PREAMBLE
  moredelim=[il][\color{red}\sout]{\%DIF\ <\ }, %DIF PREAMBLE
  moredelim=[il][\color{blue}\uwave]{\%DIF\ >\ } %DIF PREAMBLE
} %DIF PREAMBLE
\lstdefinestyle{DIFverbatimstyle}{ %DIF PREAMBLE
	language=DIFcode, %DIF PREAMBLE
	basicstyle=\ttfamily, %DIF PREAMBLE
	columns=fullflexible, %DIF PREAMBLE
	keepspaces=true %DIF PREAMBLE
} %DIF PREAMBLE
\lstnewenvironment{DIFverbatim}{\lstset{style=DIFverbatimstyle}}{} %DIF PREAMBLE
\lstnewenvironment{DIFverbatim*}{\lstset{style=DIFverbatimstyle,showspaces=true}}{} %DIF PREAMBLE
%DIF END PREAMBLE EXTENSION ADDED BY LATEXDIFF

\begin{document}

\maketitle
\DIFdelbegin %DIFDELCMD < 

%DIFDELCMD < %%%
\DIFdelend \section*{Outline and motivation}

In the previous lecture we discussed delay time tomography, this being the first and simplest tomographic method developed. 
While this approach is still used, it has been superseded by \textbf{waveform tomography}. Delay time tomography is based on the 
measurement of a small number of delay times from each seismogram, with these data modelled using linearised ray theory. 
By contrast, waveform tomography uses more of the data contained in the seismograms, fully accounts for the complexities 
of elastic wave propagation in the Earth, and does not linearise the inverse problem. This lecture will outline the ideas 
behind waveform tomography, focusing on the solution of the inverse problem using \textbf{gradient-based optimisation}.

\section*{Numerical wavefield simulations}

Waveform tomography depends on the ability to solve the full elastic wave equation numerically within 
realistic earth models. This has been possible routinely on a global scale for around 20 years, with this process 
facilitated by both algorithmic and hardware advances. These lectures are not about numerical methods. But it will 
be seen that within this subject computational costs must be carefully considered. As a number to bear in mind, the 
simulation of elastic wave propagation on a global scale within a realistic 3D earth model takes of order ten minutes
running on a modest GPU cluster. Within a tomographic study, it would be typical to use 
data from around \DIFdelbegin \DIFdel{$10^{2}-10^{3}$ earthquakes }\DIFdelend \DIFaddbegin \DIFadd{$10^{3}$ earthquakes recorded at thousands of stations (see Fig.~\ref{fig:1})}\DIFaddend , and a separate numerical 
simulation is required for each \DIFdelbegin \DIFdel{one}\DIFdelend \DIFaddbegin \DIFadd{earthquake}\DIFaddend . This means that the time for solving the necessary forward \DIFdelbegin \DIFdel{problem }\DIFdelend \DIFaddbegin \DIFadd{problems }\DIFaddend is of 
order hours or more. 

\DIFaddbegin \begin{center}
        \includegraphics[width=0.66\textwidth]{figures/L8F1.pdf}    
    \captionof{figure}{The 2160 earthquakes used to construct the global waveform tomographic model GLAD-M35 (see Fig.~\ref{fig:3}): 
    (a) locations, each event being drawn as a visual representation of its moment tensor, coloured by depth; (b) the distribution of magnitudes; (c) the distribution of 
    hypocentre depths; (d) the number of events by year. From Cui et al.\ (2024).}
    \label{fig:1}
\end{center}

\DIFaddend \section*{Waveform measurements}

Let \DIFdelbegin \DIFdel{$\mathbf{u}^{obs}(t)$ }\DIFdelend \DIFaddbegin \DIFadd{$\mathbf{u}^{\mathrm{obs}}(t)$ }\DIFaddend denote the displacement vector recorded by a seismograph at the surface location,
 $\mathbf{x}_{r}$\DIFaddbegin \DIFadd{, }\DIFaddend following an earthquake. In reality, seismographs do not record the displacement vector directly,
  with the signal instead being modified due to the mechanical and electronic properties of the instrument along with 
  additional processing done by the seismologist. For simplicity, however, we neglect such effects, but they can be 
  incorporated at the cost of only more complicated notation.

\DIFdelbegin %DIFDELCMD < \begin{figure}[ht]
%DIFDELCMD <     %%%
\DIFdelendFL \DIFaddbeginFL \begin{figure}[htbp]
    \DIFaddendFL \centering
    \DIFdelbeginFL %DIFDELCMD < \includegraphics[width=0.7\textwidth]{figures/L8F1.png}    
%DIFDELCMD <     %%%
\DIFdelendFL \DIFaddbeginFL \includegraphics[width=0.7\textwidth]{figures/L8F2.png}    
    \DIFaddendFL \caption{An example of observed and synthetic three-component seismograms from a waveform tomographic study of 
    Southern California. Data \DIFdelbeginFL \DIFdelFL{is }\DIFdelendFL \DIFaddbeginFL \DIFaddFL{are }\DIFaddendFL shown in black, and synthetics in red. On the left the comparison is made prior 
    to inversion, with the synthetics calculated in a laterally homogeneous model. The \DIFdelbeginFL \DIFdelFL{right hand }\DIFdelendFL \DIFaddbeginFL \DIFaddFL{right-hand }\DIFaddendFL side then shows the
     data fit following inversion involving 16 iterative model updates. Taken from Tape et al.\DIFaddbeginFL \DIFaddFL{\ }\DIFaddendFL (2010).}
    \DIFdelbeginFL %DIFDELCMD < \label{fig:1}
%DIFDELCMD < %%%
\DIFdelendFL \DIFaddbeginFL \label{fig:2}
\DIFaddendFL \end{figure}

Suppose that we have models for both the Earth and earthquake, and so can numerically solve the elastic wave equation to 
obtain a synthetic seismogram $\mathbf{u}(\mathbf{x}_{r},t)$ at this location. To quantify the misfit between the observed and
 synthetic seismograms we can define the \DIFdelbegin \DIFdel{least squares misfit
}\begin{displaymath}
\DIFdel{J=\frac{1}{2}\int_{0}^{T}||\mathbf{u}(\mathbf{x}_{r},t)-\mathbf{u}^{obs}(t)||^{2}dt,
}\end{displaymath}%DIFAUXCMD
\DIFdelend \DIFaddbegin \DIFadd{least-squares misfit
}\begin{equation}
\DIFadd{J=\frac{1}{2}\int_{0}^{T}\|\mathbf{u}(\mathbf{x}_{r},t)-\mathbf{u}^{\mathrm{obs}}(t)\|^{2}\dd t,
\label{eq:1}
}\end{equation}\DIFaddend 
with $[0,T]$ being the time interval over which measurements have been made. Recorded seismograms are, of course, discretely 
sampled, and so really we should write a sum and not an 
integral over time. But in practice sampling rates are sufficiently high that this distinction can be ignored.

 The value of $J$ depends on the synthetic seismogram, $\mathbf{u}$, and hence ultimately on the underlying 
 parameters for the earth model and the earthquake. We could, therefore, try to minimise $J$ in the 
 hope that the model parameters are then closer \DIFaddbegin \DIFadd{to }\DIFaddend their true values. In practice, of course, we would 
 not actually use $J$ as defined above, but sum such terms over many different locations and earthquakes,
and regularisation terms might also be added. \DIFaddbegin \DIFadd{Fig.~\ref{fig:2} shows the sort of improvement in the fit between observed 
and synthetic seismograms that can be achieved by such an inversion. }\DIFaddend In fact, \DIFaddbegin \DIFadd{the }\DIFaddend least-squares misfit between seismograms is rarely 
used in practice, with other more complicated measures of misfit being preferred. Nevertheless, 
this simple functional is sufficient to motivate and illustrate \DIFdelbegin \DIFdel{wave form }\DIFdelend \DIFaddbegin \DIFadd{waveform }\DIFaddend inversion.

\section*{Functional derivatives and sensitivity kernels}
The waveform misfit, $J$, defined above depends on the synthetic wavefield, $\mathbf{u}$, and hence we can write it
 as $J(\mathbf{u})$. But the value of $\mathbf{u}$ is determined through solution of the equations of motion from the 
 underlying model parameters. To simplify \DIFdelbegin \DIFdel{notations}\DIFdelend \DIFaddbegin \DIFadd{notation}\DIFaddend , let us write $\mathbf{m}$ for these model parameters. Here $\mathbf{m}$ could 
 be all of $\rho$, $A_{ijkl}$ and $\overline{S}_{ij}$, but in practice it is often some subset of these parameters. For example, 
 within a tomographic problem it might be assumed that the source parameters are known. 

 For a given value of $\mathbf{m}$ we can 
 compute the corresponding synthetic wavefield, $\mathbf{u}$, and from this the value of the misfit, $J(\mathbf{u})$. To indicate that 
 the misfit can be viewed as an implicit function of $\mathbf{m}$, we will write it as $J(\mathbf{m})$. This is a slight abuse of notation, 
 but it should always be clear from context which variables the misfit depends on.

Within waveform tomography, the goal is to find model parameters that explain the observed data, and this is done by 
minimising $\mathbf{m}\mapsto J(\mathbf{m})$. For a small model perturbation, $\delta \mathbf{m}$, we can write to first-order accuracy
\DIFdelbegin \begin{displaymath}
\DIFdel{J(\mathbf{m}+\delta \mathbf{m})=J(\mathbf{m})+\langle DJ(\mathbf{m})|\delta \mathbf{m}\rangle+O(||\delta \mathbf{m}||^{2}),
}\end{displaymath}%DIFAUXCMD
\DIFdelend \DIFaddbegin \begin{equation}
\DIFadd{J(\mathbf{m}+\delta \mathbf{m})=J(\mathbf{m})+\langle DJ(\mathbf{m})|\delta \mathbf{m}\rangle+O(\|\delta \mathbf{m}\|^{2}),
\label{eq:2}
}\end{equation}\DIFaddend 
where $DJ(\mathbf{m})$ denotes the functional derivative of \DIFdelbegin \DIFdel{J }\DIFdelend \DIFaddbegin \DIFadd{$J$ }\DIFaddend with respect to $\mathbf{m}$ and $\langle \cdot | \cdot \rangle$ is an 
appropriate inner product. As an example, supposing that the shear velocity, $\beta$, is the only model parameter\DIFaddbegin \DIFadd{, }\DIFaddend we can concretely write
\DIFdelbegin \begin{displaymath}
\DIFdel{\langle DJ(\mathbf{m})|\delta \mathbf{m}\rangle=\int_{M}K_{\beta}\delta\beta\dd^{3}\mathbf{x}
}\end{displaymath}%DIFAUXCMD
\DIFdelend \DIFaddbegin \begin{equation}
\DIFadd{\langle DJ(\mathbf{m})|\delta \mathbf{m}\rangle=\int_{M}K_{\beta}\,\delta\beta\dd^{3}\mathbf{x},
\label{eq:3}
}\end{equation}\DIFaddend 
where the function $K_{\beta}$ is known as a \textbf{sensitivity kernel} for $J$ with respect to $\beta$; 
note that in this context the terms functional derivative, gradient, and sensitivity kernel are used interchangeably.
We discuss later how these functional derivatives can be obtained practically, but for the moment we just assume that this can be done.

\section*{\DIFdelbegin \DIFdel{Gradient based }\DIFdelend \DIFaddbegin \DIFadd{Gradient-based }\DIFaddend optimisation}
If $\overline{\mathbf{m}}$ is a minimum of $J$, then for any sufficiently small $\delta \mathbf{m}$ we should have
\DIFdelbegin \begin{displaymath}
\DIFdel{J(\overline{\mathbf{m}}+\delta \mathbf{m})\ge J(\overline{\mathbf{m}}),
}\end{displaymath}%DIFAUXCMD
\DIFdelend \DIFaddbegin \begin{equation}
\DIFadd{J(\overline{\mathbf{m}}+\delta \mathbf{m})\ge J(\overline{\mathbf{m}}),
\label{eq:4}
}\end{equation}\DIFaddend 
which is true if and only if
\DIFdelbegin \begin{displaymath}
\DIFdel{DJ(\overline{\mathbf{m}})=\mathbf{0}.
}\end{displaymath}%DIFAUXCMD
\DIFdelend \DIFaddbegin \begin{equation}
\DIFadd{DJ(\overline{\mathbf{m}})=\mathbf{0}.
\label{eq:5}
}\end{equation}\DIFaddend 
This condition is, of course, familiar from calculus in finite-dimensional spaces. 

There are a number of approaches to solve this optimisation problem. One is a systematic \DIFdelbegin \DIFdel{"grid search" }\DIFdelend \DIFaddbegin \DIFadd{``grid search'' }\DIFaddend of the model space. 
Doing this would require the model space to be discretised using a finite-dimensional basis as in Lecture\DIFdelbegin \DIFdel{18. }\DIFdelend \DIFaddbegin \DIFadd{~19. }\DIFaddend The issue with 
this method is the computational cost. Suppose, for example, that the discretised model space had \DIFdelbegin \DIFdel{n-dimensions}\DIFdelend \DIFaddbegin \DIFadd{$n$ dimensions}\DIFaddend , and that along 
each dimension experience had shown that \DIFdelbegin \DIFdel{m }\DIFdelend \DIFaddbegin \DIFadd{$m$ }\DIFaddend points needed to be tested. The number of forward calculations required would then be
$m^{n}$. In realistic problems, \DIFdelbegin \DIFdel{n }\DIFdelend \DIFaddbegin \DIFadd{$n$ }\DIFaddend is of order thousands or millions, while we have already noted the high cost of solving the 
forward problem. The result is that grid search methods are not viable. 

Random sampling methods such as Markov \DIFdelbegin \DIFdel{Chain }\DIFdelend \DIFaddbegin \DIFadd{chain }\DIFaddend Monte Carlo (MCMC) simulations perform better within high-dimensional spaces 
(the number of forward calculations scales with something like $n^{2}$), but the costs are still far too high to be used 
practically in waveform tomography. 

As a result, tomographers must rely on local optimisation methods that seek a point such that eq.\DIFdelbegin \DIFdel{(5}\DIFdelend \DIFaddbegin \DIFadd{~(\ref{eq:5}}\DIFaddend ) holds. 
An obvious downside to this approach is that the model obtained may only be a \textbf{local and not a global minimum of 
the misfit}. But this is typically all that is currently possible, and there do exist approximate but useful methods available 
for uncertainty quantification\DIFdelbegin \footnote{\DIFdel{About which there is not sufficient time to discuss.}}%DIFAUXCMD
\addtocounter{footnote}{-1}%DIFAUXCMD
\DIFdel{.}\DIFdelend \DIFaddbegin \DIFadd{.}\footnote{\DIFadd{About which there is not sufficient time to discuss.}}
\DIFaddend 

Suppose we start from an initial model $\mathbf{m}_{0}$. At this point we can calculate the value of the misfit,
 $J(\mathbf{m}_{0})$, along with its gradient, $DJ(\mathbf{m}_{0})$. The gradient points in the direction along which 
 the misfit most rapidly increases, and the negative gradient gives the direction of most rapid decrease. This suggests
  that we can look for a new model of the form
\DIFdelbegin \begin{displaymath}
\DIFdel{\mathbf{m}_{1}=\mathbf{m}_{0}-\lambda~DJ(\mathbf{m}_{0})
}\end{displaymath}%DIFAUXCMD
\DIFdelend \DIFaddbegin \begin{equation}
\DIFadd{\mathbf{m}_{1}=\mathbf{m}_{0}-\lambda\,DJ(\mathbf{m}_{0}),
\label{eq:6}
}\end{equation}\DIFaddend 
where $\lambda>0$ is some \DIFdelbegin \DIFdel{step-length }\DIFdelend \DIFaddbegin \DIFadd{step length }\DIFaddend to be determined. To find a suitable value for $\lambda$, we need only consider 
a one-dimensional minimisation problem for the function
\DIFdelbegin \begin{displaymath}
\DIFdel{f(\lambda)=J}[\DIFdel{\mathbf{m}_{0}-\lambda~DJ(\mathbf{m}_{0})}]\DIFdel{,
}\end{displaymath}%DIFAUXCMD
\DIFdelend \DIFaddbegin \begin{equation}
\DIFadd{f(\lambda)=J}[\DIFadd{\mathbf{m}_{0}-\lambda\,DJ(\mathbf{m}_{0})}]\DIFadd{,
\label{eq:7}
}\end{equation}\DIFaddend 
which can be solved quite easily. Having now arrived at a new model, $\mathbf{m}_{1}$, with a lower 
misfit\DIFaddbegin \DIFadd{, }\DIFaddend we can simply repeat the whole process until the gradient vanishes and a local minimum has been 
found\DIFdelbegin \footnote{\DIFdel{In practice, the process might be stopped when the misfit has decreased sufficiently much from a
 statistical perspective.}}%DIFAUXCMD
\addtocounter{footnote}{-1}%DIFAUXCMD
\DIFdel{.}\DIFdelend \DIFaddbegin \DIFadd{.}\footnote{\DIFadd{In practice, the process might be stopped when the misfit has decreased sufficiently from a
 statistical perspective.}} \DIFaddend The above method is the simplest example of gradient-based optimisation\DIFaddbegin \DIFadd{, and is }\DIFaddend known as
\textbf{the method of steepest descent}. More sophisticated algorithms for choosing the descent direction\DIFaddbegin \DIFadd{, 
which have superior convergence properties, }\DIFaddend tend to be used in practice\DIFdelbegin \DIFdel{which have superior convergence properties}\DIFdelend . In all cases, however, this direction 
is built from linear combinations of the current and past gradients of the misfit.

\section*{How not to calculate the gradient of $J$}
\DIFdelbegin \DIFdel{Gradient based }\DIFdelend \DIFaddbegin \DIFadd{Gradient-based }\DIFaddend optimisation depends on being able \DIFdelbegin \DIFdel{to both }\DIFdelend \DIFaddbegin \DIFadd{both to }\DIFaddend solve the forward problem \DIFdelbegin \DIFdel{, }\DIFdelend and to 
calculate the gradient of the misfit. As noted, the first task can be done using numerical wavefield simulations 
at a high but manageable cost. Schematically, such a code provides a means for performing the 
mapping $\mathbf{m}\mapsto J(\mathbf{m})$. With this at hand, one way to get the required gradient is to make use of 
the \textbf{finite-difference approximation}
\DIFdelbegin \begin{displaymath}
\DIFdel{\langle DJ(\mathbf{m})|\delta \mathbf{m}\rangle\approx J(\mathbf{m}+\delta \mathbf{m})-J(\mathbf{m}).
}\end{displaymath}%DIFAUXCMD
\DIFdelend \DIFaddbegin \begin{equation}
\DIFadd{\langle DJ(\mathbf{m})|\delta \mathbf{m}\rangle\approx J(\mathbf{m}+\delta \mathbf{m})-J(\mathbf{m}).
\label{eq:8}
}\end{equation}\DIFaddend 
Using this formula, we can get the derivative of the misfit along any direction at the cost of two forward calculations.
Suppose we introduce an $m$-dimensional
parameterisation for the model space. Using $m+1$ forward calculations we could then approximate the gradient relative to this
basis. The advantage of this approach is that it requires no additional theory or coding once a numerical wave propagation 
code has been developed. The problem is, as noted before, \DIFaddbegin \DIFadd{that }\DIFaddend the number of model parameters, $m$, is typically large\DIFaddbegin \DIFadd{, }\DIFaddend and so the
overall cost is too high to compute a single gradient in this manner, let alone the many different gradients needed in \DIFaddbegin \DIFadd{the 
}\DIFaddend iterative solution of the optimisation problem.

\section*{How you should calculate the gradient of $J$}
Our aim is to calculate the gradient of $J(\mathbf{m})$. This is equivalent to finding the gradient 
of $J(\mathbf{u})$ subject to the constraint that $\mathbf{u}$ is related to $\mathbf{m}$ through solution 
of the forward problem. To proceed, we make use of the method of \textbf{Lagrange multipliers}\DIFaddbegin \DIFadd{, }\DIFaddend which many 
of you will have seen previously, though not necessarily applied to fields. Within this course it is necessary 
to know how to apply this method, but not to explain why it works. For those interested\DIFaddbegin \DIFadd{, }\DIFaddend a non-examinable 
outline is provided in the appendix to this lecture. We begin by defining a \textbf{Lagrangian} for the problem\DIFaddbegin \DIFadd{, }\DIFaddend which takes the following form
\DIFdelbegin \begin{align*}
\DIFdel{L(\mathbf{u},\mathbf{u}',\mathbf{w}',\mathbf{m}) = J(\mathbf{u}) }&\DIFdel{- \int_{0}^{T}\int_{M}\left[\rho\frac{\partial^{2}u_{i}}{\partial t^{2}}-\frac{\partial}{\partial x_{j}}\left(A_{ijkl}\frac{\partial u_{k}}{\partial x_{l}}-\overline{S}_{ij}\right)\right]u_{i}'\dd^{3}\mathbf{x}\dd t \nonumber }\\
&\DIFdel{- \int_{0}^{T}\int_{\partial M}\left(A_{ijkl}\frac{\partial u_{k}}{\partial x_{l}}-\overline{S}_{ij}\right)\hat{n}_{j}w_{i}'\dd S\dd t.
}\end{align*}%DIFAUXCMD
\DIFdelend \DIFaddbegin \begin{align}
\DIFadd{L(\mathbf{u},\mathbf{u}',\mathbf{w}',\mathbf{m}) = J(\mathbf{u}) }&\DIFadd{- \int_{0}^{T}\int_{M}\left[\rho\frac{\partial^{2}u_{i}}{\partial t^{2}}-\frac{\partial}{\partial x_{j}}\left(A_{ijkl}\frac{\partial u_{k}}{\partial x_{l}}-\overline{S}_{ij}\right)\right]u_{i}'\dd^{3}\mathbf{x}\dd t \nonumber }\\
&\DIFadd{- \int_{0}^{T}\int_{\partial M}\left(A_{ijkl}\frac{\partial u_{k}}{\partial x_{l}}-\overline{S}_{ij}\right)\hat{n}_{j}w_{i}'\dd S\dd t.
\label{eq:9}
}\end{align}\DIFaddend 
Here we see first on the \DIFdelbegin \DIFdel{right hand }\DIFdelend \DIFaddbegin \DIFadd{right-hand }\DIFaddend side the unconstrained misfit, $J(\mathbf{u})$. The second term is associated with 
the constraint that $\mathbf{u}$ should satisfy the equations of motion and includes a 
Lagrange multiplier field $\mathbf{u}'$. The final term does the same for the boundary conditions, and 
features a second \DIFdelbegin \DIFdel{Lagrangian }\DIFdelend \DIFaddbegin \DIFadd{Lagrange }\DIFaddend multiplier field $\mathbf{w}'$. 

The \textbf{Lagrange multiplier theorem} tells us that the equality
\DIFdelbegin \begin{displaymath}
\DIFdel{DJ(\mathbf{m})=D_{\mathbf{m}}L(\mathbf{u},\mathbf{u}',\mathbf{w}',\mathbf{m}),
}\end{displaymath}%DIFAUXCMD
\DIFdelend \DIFaddbegin \begin{equation}
\DIFadd{DJ(\mathbf{m})=D_{\mathbf{m}}L(\mathbf{u},\mathbf{u}',\mathbf{w}',\mathbf{m})
\label{eq:10}
}\end{equation}\DIFaddend 
holds so long as the following conditions are met
\begin{align}
D_{\mathbf{u}}L(\mathbf{u},\mathbf{u}',\mathbf{w}',\mathbf{m})&=0, \label{eq:11} \\
D_{\mathbf{u}'}L(\mathbf{u},\mathbf{u}',\mathbf{w}',\mathbf{m})&=0, \label{eq:12} \\
D_{\mathbf{w}'}L(\mathbf{u},\mathbf{u}',\mathbf{w}',\mathbf{m})&=0.
\label{eq:13}
\end{align}
Here subscripts are used for partial functional derivatives in an obvious manner. To understand this result, 
we begin with eq.\DIFdelbegin \DIFdel{\eqref{eq:12}}\DIFdelend \DIFaddbegin \DIFadd{~(\ref{eq:12})}\DIFaddend . The dependence of the Lagrangian on the field $\mathbf{u}'$ is very simple, 
and so we readily find that
\DIFdelbegin \begin{displaymath}
\DIFdel{\langle D_{\mathbf{u}'}L(\mathbf{u},\mathbf{u}',\mathbf{w}',\mathbf{m})|\delta \mathbf{u}'\rangle = -\int_{0}^{T}\int_{M}\left[\rho\frac{\partial^{2}u_{i}}{\partial t^{2}}-\frac{\partial}{\partial x_{j}}\left(A_{ijkl}
\frac{\partial u_{k}}{\partial x_{l}}-\overline{S}_{ij}\right)\right]\delta u'_{i}\dd^{3}\mathbf{x}dt.
}\end{displaymath}%DIFAUXCMD
\DIFdelend \DIFaddbegin \begin{equation}
\DIFadd{\langle D_{\mathbf{u}'}L(\mathbf{u},\mathbf{u}',\mathbf{w}',\mathbf{m})|\delta \mathbf{u}'\rangle = -\int_{0}^{T}\int_{M}\left[\rho\frac{\partial^{2}u_{i}}{\partial t^{2}}-\frac{\partial}{\partial x_{j}}\left(A_{ijkl}
\frac{\partial u_{k}}{\partial x_{l}}-\overline{S}_{ij}\right)\right]\delta u'_{i}\dd^{3}\mathbf{x}\dd t.
\label{eq:14}
}\end{equation}\DIFaddend 
We need this expression to vanish for all possible choices of $\delta \mathbf{u}'$, and hence we 
conclude that $\mathbf{u}$ must satisfy the expected equations of motion
\DIFdelbegin \begin{displaymath}
\DIFdel{\rho\frac{\partial^{2}u_{i}}{\partial t^{2}}-\frac{\partial}{\partial x_{j}}\left(A_{ijkl}\frac{\partial u_{k}}{\partial x_{l}}-\overline{S}_{ij}\right)=0.
}\end{displaymath}%DIFAUXCMD
\DIFdelend \DIFaddbegin \begin{equation}
\DIFadd{\rho\frac{\partial^{2}u_{i}}{\partial t^{2}}-\frac{\partial}{\partial x_{j}}\left(A_{ijkl}\frac{\partial u_{k}}{\partial x_{l}}-\overline{S}_{ij}\right)=0.
\label{eq:15}
}\end{equation}\DIFaddend 
Similarly, the final condition in eq.\DIFdelbegin \DIFdel{\eqref{eq:13} }\DIFdelend \DIFaddbegin \DIFadd{~(\ref{eq:13}) }\DIFaddend gives the boundary conditions for $\mathbf{u}$. 

Turning now to the first condition in eq.\DIFdelbegin \DIFdel{\eqref{eq:11}}\DIFdelend \DIFaddbegin \DIFadd{~(\ref{eq:11})}\DIFaddend , we find that for any $\delta \mathbf{u}$
\DIFdelbegin \begin{align*}
\DIFdel{\langle D_{\mathbf{u}}L(\mathbf{u},\mathbf{u}',\mathbf{w}',\mathbf{m})|\delta \mathbf{u}\rangle = \langle DJ(\mathbf{u})|\delta \mathbf{u}\rangle }&\DIFdel{- \int_{0}^{T}\int_{M}\left[\rho\frac{\partial^{2}\delta u_{i}}{\partial t^{2}}-\frac{\partial}{\partial x_{j}}
\left(A_{ijkl}\frac{\partial\delta u_{k}}{\partial x_{l}}\right)\right]u_{i}'\dd^{3}\mathbf{x}dt \nonumber }\\
&\DIFdel{- \int_{0}^{T}\int_{\partial M}\left(A_{ijkl}\frac{\partial\delta u_{k}}{\partial x_{l}}\right)\hat{n}_{j}w_{i}'\dd S\dd t,
}\end{align*}%DIFAUXCMD
\DIFdelend \DIFaddbegin \begin{align}
\DIFadd{\langle D_{\mathbf{u}}L(\mathbf{u},\mathbf{u}',\mathbf{w}',\mathbf{m})|\delta \mathbf{u}\rangle = \langle DJ(\mathbf{u})|\delta \mathbf{u}\rangle }&\DIFadd{- \int_{0}^{T}\int_{M}\left[\rho\frac{\partial^{2}\delta u_{i}}{\partial t^{2}}-\frac{\partial}{\partial x_{j}}
\left(A_{ijkl}\frac{\partial\delta u_{k}}{\partial x_{l}}\right)\right]u_{i}'\dd^{3}\mathbf{x}\dd t \nonumber }\\
&\DIFadd{- \int_{0}^{T}\int_{\partial M}\left(A_{ijkl}\frac{\partial\delta u_{k}}{\partial x_{l}}\right)\hat{n}_{j}w_{i}'\dd S\dd t,
\label{eq:16}
}\end{align}\DIFaddend 
must vanish. This is a more complicated expression, and understanding its consequences will take some effort. 
But we know that $\mathbf{u}$ is fixed uniquely by the previous constraints, and hence the above expression 
must somehow determine the Lagrange multiplier fields. As a first step, we can use eq.\DIFdelbegin \DIFdel{(1}\DIFdelend \DIFaddbegin \DIFadd{~(\ref{eq:1}}\DIFaddend ) to explicitly write 
the first term on the \DIFdelbegin \DIFdel{right hand side as
}\begin{displaymath}
\DIFdel{\langle DJ(\mathbf{u})|\delta \mathbf{u}\rangle=\int_{0}^{T}[u_{i}(\mathbf{x}_{r},t)-u_{i}^{obs}(t)]\delta u_{i}(\mathbf{x}_{r},t)dt.
}\end{displaymath}%DIFAUXCMD
\DIFdelend \DIFaddbegin \DIFadd{right-hand side as
}\begin{equation}
\DIFadd{\langle DJ(\mathbf{u})|\delta \mathbf{u}\rangle=\int_{0}^{T}[u_{i}(\mathbf{x}_{r},t)-u_{i}^{\mathrm{obs}}(t)]\delta u_{i}(\mathbf{x}_{r},t)\dd t.
\label{eq:17}
}\end{equation}\DIFaddend 
Making use of the Dirac delta function, this can be conveniently written
\DIFdelbegin \begin{displaymath}
\DIFdel{\langle DJ(\mathbf{u})|\delta \mathbf{u}\rangle=\int_{0}^{T}\int_{\partial M}h_{i}'\delta u_{i}\dd S\dd t,
}\end{displaymath}%DIFAUXCMD
\DIFdelend \DIFaddbegin \begin{equation}
\DIFadd{\langle DJ(\mathbf{u})|\delta \mathbf{u}\rangle=\int_{0}^{T}\int_{\partial M}h_{i}'\delta u_{i}\dd S\dd t,
\label{eq:18}
}\end{equation}\DIFaddend 
where we have set
\DIFdelbegin \begin{displaymath}
\DIFdel{h_{i}'(\mathbf{x},t)=}[\DIFdel{u_{i}(\mathbf{x}_{r},t)-u_{i}^{obs}(t)}]\DIFdel{\delta(\mathbf{x}-\mathbf{x}_{r}).
}\end{displaymath}%DIFAUXCMD
\DIFdelend \DIFaddbegin \begin{equation}
\DIFadd{h_{i}'(\mathbf{x},t)=}[\DIFadd{u_{i}(\mathbf{x}_{r},t)-u_{i}^{\mathrm{obs}}(t)}]\DIFadd{\delta(\mathbf{x}-\mathbf{x}_{r}).
\label{eq:19}
}\end{equation}\DIFaddend 
With this done, we can write the required condition as
\DIFdelbegin \begin{align*}
\DIFdel{\int_{0}^{T}\int_{\partial M}h_{i}'\delta u_{i}\dd S\dd t }&\DIFdel{- \int_{0}^{T}\int_{M}\left[\rho\frac{\partial^{2}\delta u_{i}}{\partial t^{2}}-\frac{\partial}{\partial x_{j}}
\left(A_{ijkl}\frac{\partial\delta u_{k}}{\partial x_{l}}\right)\right]u_{i}'\dd^{3}\mathbf{x}\dd t \nonumber }\\
&\DIFdel{- \int_{0}^{T}\int_{\partial M}\left(A_{ijkl}\frac{\partial\delta u_{k}}{\partial x_{l}}\right)\hat{n}_{j}w_{i}'\dd S\dd t=0,
}\end{align*}%DIFAUXCMD
\DIFdelend \DIFaddbegin \begin{align}
\DIFadd{\int_{0}^{T}\int_{\partial M}h_{i}'\delta u_{i}\dd S\dd t }&\DIFadd{- \int_{0}^{T}\int_{M}\left[\rho\frac{\partial^{2}\delta u_{i}}{\partial t^{2}}-\frac{\partial}{\partial x_{j}}
\left(A_{ijkl}\frac{\partial\delta u_{k}}{\partial x_{l}}\right)\right]u_{i}'\dd^{3}\mathbf{x}\dd t \nonumber }\\
&\DIFadd{- \int_{0}^{T}\int_{\partial M}\left(A_{ijkl}\frac{\partial\delta u_{k}}{\partial x_{l}}\right)\hat{n}_{j}w_{i}'\dd S\dd t=0,
\label{eq:20}
}\end{align}\DIFaddend 
for all $\delta \mathbf{u}$. 

To proceed, we need to successively integrate by parts to shift derivatives from the perturbed 
displacement field to the Lagrange multipliers. Starting with the \DIFdelbegin \DIFdel{time-derivatives}\DIFdelend \DIFaddbegin \DIFadd{time derivatives}\DIFaddend , we have
\DIFdelbegin \begin{displaymath}
\DIFdel{\int_{0}^{T}\frac{\partial^{2}\delta u_{i}}{\partial t^{2}}u_{i}'dt = -\int_{0}^{T}\frac{\partial\delta u_{i}}{\partial t}\frac{\partial u_{i}'}{\partial t}dt+\left[\frac{\partial\delta u_{i}}{\partial t}u_{i}'\right]_{t=T},
}\end{displaymath}%DIFAUXCMD
\DIFdelend \DIFaddbegin \begin{equation}
\DIFadd{\int_{0}^{T}\frac{\partial^{2}\delta u_{i}}{\partial t^{2}}u_{i}'\dd t = -\int_{0}^{T}\frac{\partial\delta u_{i}}{\partial t}\frac{\partial u_{i}'}{\partial t}\dd t+\left[\frac{\partial\delta u_{i}}{\partial t}u_{i}'\right]_{t=T},
\label{eq:21}
}\end{equation}\DIFaddend 
where we have used the vanishing initial conditions on $u_{i}$ (and hence $\delta u_{i}$) to eliminate one 
of the boundary terms. Integrating by parts once more gives
\DIFdelbegin \begin{displaymath}
\DIFdel{\int_{0}^{T}\frac{\partial^{2}\delta u_{i}}{\partial t^{2}}u_{i}'dt = \int_{0}^{T}\delta u_{i}\frac{\partial^{2}u_{i}'}{\partial t^{2}}dt+\left[\frac{\partial\delta u_{i}}{\partial t}u_{i}'\right]_{t=T}-\left[\delta u_{i}\frac{\partial u_{i}'}{\partial t}\right]_{t=T}.
}\end{displaymath}%DIFAUXCMD
\DIFdelend \DIFaddbegin \begin{equation}
\DIFadd{\int_{0}^{T}\frac{\partial^{2}\delta u_{i}}{\partial t^{2}}u_{i}'\dd t = \int_{0}^{T}\delta u_{i}\frac{\partial^{2}u_{i}'}{\partial t^{2}}\dd t+\left[\frac{\partial\delta u_{i}}{\partial t}u_{i}'\right]_{t=T}-\left[\delta u_{i}\frac{\partial u_{i}'}{\partial t}\right]_{t=T}.
\label{eq:22}
}\end{equation}\DIFaddend 
Looking now at the spatial derivatives, we first obtain
\DIFdelbegin \begin{align*}
\DIFdel{\int_{M}\frac{\partial}{\partial x_{j}}}&\DIFdel{\left(A_{ijkl}\frac{\partial\delta u_{k}}{\partial x_{l}}\right)u_{i}'\dd^{3}\mathbf{x} - \int_{\partial M}\left(A_{ijkl}\frac{\partial\delta u_{k}}{\partial x_{l}}\right)\hat{n}_{j}w_{i}'\dd S \nonumber }\\
&\DIFdel{= -\int_{M}A_{ijkl}\frac{\partial\delta u_{k}}{\partial x_{l}}\frac{\partial u_{i}'}{\partial x_{j}}\dd^{3}\mathbf{x} - \int_{\partial M}\left(A_{ijkl}\frac{\partial\delta u_{k}}{\partial x_{l}}\right)\hat{n}_{j}(w_{i}'-u_{i}')\dd S.
}\end{align*}%DIFAUXCMD
\DIFdelend \DIFaddbegin \begin{align}
\DIFadd{\int_{M}\frac{\partial}{\partial x_{j}}}&\DIFadd{\left(A_{ijkl}\frac{\partial\delta u_{k}}{\partial x_{l}}\right)u_{i}'\dd^{3}\mathbf{x} - \int_{\partial M}\left(A_{ijkl}\frac{\partial\delta u_{k}}{\partial x_{l}}\right)\hat{n}_{j}w_{i}'\dd S \nonumber }\\
&\DIFadd{= -\int_{M}A_{ijkl}\frac{\partial\delta u_{k}}{\partial x_{l}}\frac{\partial u_{i}'}{\partial x_{j}}\dd^{3}\mathbf{x} - \int_{\partial M}\left(A_{ijkl}\frac{\partial\delta u_{k}}{\partial x_{l}}\right)\hat{n}_{j}(w_{i}'-u_{i}')\dd S.
\label{eq:23}
}\end{align}\DIFaddend 
Making use of the hyperelastic symmetry, $A_{ijkl}=A_{klij}$\DIFaddbegin \DIFadd{, }\DIFaddend we can then integrate by parts once more to arrive at
\DIFdelbegin \begin{align*}
\DIFdel{\int_{M}\frac{\partial}{\partial x_{j}}}&\DIFdel{\left(A_{ijkl}\frac{\partial\delta u_{k}}{\partial x_{l}}\right)u_{i}'\dd^{3}\mathbf{x} - \int_{\partial M}\left(A_{ijkl}\frac{\partial\delta u_{k}}{\partial x_{l}}\right)\hat{n}_{j}w_{i}'\dd S \nonumber }\\
&\DIFdel{= \int_{M}\frac{\partial}{\partial x_{j}}\left(A_{ijkl}\frac{\partial u_{k}'}{\partial x_{l}}\right)\delta u_{i}\dd^{3}\mathbf{x} - \int_{\partial M}\left(A_{ijkl}\frac{\partial\delta u_{k}}{\partial x_{l}}\right)\hat{n}_{j}(w_{i}'-u_{i}')\dd S \nonumber }\\
& \DIFdel{\qquad -\int_{\partial M}\left(A_{ijkl}\frac{\partial u_{k}'}{\partial x_{l}}\right)\hat{n}_{j}\delta u_{i}\dd S.
}\end{align*}%DIFAUXCMD
\DIFdelend \DIFaddbegin \begin{align}
\DIFadd{\int_{M}\frac{\partial}{\partial x_{j}}}&\DIFadd{\left(A_{ijkl}\frac{\partial\delta u_{k}}{\partial x_{l}}\right)u_{i}'\dd^{3}\mathbf{x} - \int_{\partial M}\left(A_{ijkl}\frac{\partial\delta u_{k}}{\partial x_{l}}\right)\hat{n}_{j}w_{i}'\dd S \nonumber }\\
&\DIFadd{= \int_{M}\frac{\partial}{\partial x_{j}}\left(A_{ijkl}\frac{\partial u_{k}'}{\partial x_{l}}\right)\delta u_{i}\dd^{3}\mathbf{x} - \int_{\partial M}\left(A_{ijkl}\frac{\partial\delta u_{k}}{\partial x_{l}}\right)\hat{n}_{j}(w_{i}'-u_{i}')\dd S \nonumber }\\
& \DIFadd{\qquad -\int_{\partial M}\left(A_{ijkl}\frac{\partial u_{k}'}{\partial x_{l}}\right)\hat{n}_{j}\delta u_{i}\dd S.
\label{eq:24}
}\end{align}\DIFaddend 



Putting these results together, we have shown that eq.\DIFdelbegin \DIFdel{(20}\DIFdelend \DIFaddbegin \DIFadd{~(\ref{eq:20}}\DIFaddend ) is equivalent to
\DIFdelbegin \begin{align*}
\DIFdel{-}&\DIFdel{\int_{0}^{T}\int_{M}\left[\rho\frac{\partial^{2}u_{i}'}{\partial t^{2}}-\frac{\partial}{\partial x_{j}}\left(A_{ijkl}\frac{\partial u_{k}'}{\partial x_{l}}\right)\right]\delta u_{i}\dd^{3}\mathbf{x}\dd t \nonumber }\\
\DIFdel{-}&\DIFdel{\int_{0}^{T}\int_{\partial M}\left(A_{ijkl}\frac{\partial\delta u_{k}}{\partial x_{l}}\right)\hat{n}_{j}(w_{i}'-u_{i}')\dd S\dd t \nonumber }\\
\DIFdel{+}&\DIFdel{\int_{0}^{T}\int_{\partial M}\left[h_{i}'-\left(A_{ijkl}\frac{\partial u_{k}'}{\partial x_{l}}\right)\hat{n}_{j}\right]\delta u_{i}\dd S\dd t \nonumber }\\
\DIFdel{-}&\DIFdel{\int_{M}\rho\left\{\left[\frac{\partial\delta u_{i}}{\partial t}u_{i}'\right]_{t=T}-\left[\delta u_{i}\frac{\partial u_{i}'}{\partial t}\right]_{t=T}\right\}\dd^{3}\mathbf{x}=0,
}\end{align*}%DIFAUXCMD
\DIFdelend \DIFaddbegin \begin{align}
\DIFadd{-}&\DIFadd{\int_{0}^{T}\int_{M}\left[\rho\frac{\partial^{2}u_{i}'}{\partial t^{2}}-\frac{\partial}{\partial x_{j}}\left(A_{ijkl}\frac{\partial u_{k}'}{\partial x_{l}}\right)\right]\delta u_{i}\dd^{3}\mathbf{x}\dd t \nonumber }\\
\DIFadd{-}&\DIFadd{\int_{0}^{T}\int_{\partial M}\left(A_{ijkl}\frac{\partial\delta u_{k}}{\partial x_{l}}\right)\hat{n}_{j}(w_{i}'-u_{i}')\dd S\dd t \nonumber }\\
\DIFadd{+}&\DIFadd{\int_{0}^{T}\int_{\partial M}\left[h_{i}'-\left(A_{ijkl}\frac{\partial u_{k}'}{\partial x_{l}}\right)\hat{n}_{j}\right]\delta u_{i}\dd S\dd t \nonumber }\\
\DIFadd{-}&\DIFadd{\int_{M}\rho\left\{\left[\frac{\partial\delta u_{i}}{\partial t}u_{i}'\right]_{t=T}-\left[\delta u_{i}\frac{\partial u_{i}'}{\partial t}\right]_{t=T}\right\}\dd^{3}\mathbf{x}=0,
\label{eq:25}
}\end{align}\DIFaddend 
for all $\delta \mathbf{u}$. From the first term in this equation, we conclude that the Lagrange 
multiplier field $\mathbf{u}'$ must satisfy the elastic wave equation
\DIFdelbegin \begin{displaymath}
\DIFdel{\rho\frac{\partial^{2}u_{i}'}{\partial t^{2}}-\frac{\partial}{\partial x_{j}}\left(A_{ijkl}\frac{\partial u_{k}'}{\partial x_{l}}\right)=0.
}\end{displaymath}%DIFAUXCMD
\DIFdelend \DIFaddbegin \begin{equation}
\DIFadd{\rho\frac{\partial^{2}u_{i}'}{\partial t^{2}}-\frac{\partial}{\partial x_{j}}\left(A_{ijkl}\frac{\partial u_{k}'}{\partial x_{l}}\right)=0.
\label{eq:26}
}\end{equation}\DIFaddend 
In dealing with the boundary terms, we need the fact that\DIFaddbegin \DIFadd{, }\DIFaddend for an arbitrary function\DIFaddbegin \DIFadd{, }\DIFaddend its boundary value and that of 
its normal derivative can be chosen separately. As a result, the second term implies that \DIFdelbegin \DIFdel{$w'$ }\DIFdelend \DIFaddbegin \DIFadd{$\mathbf{w}'$ }\DIFaddend is the restriction 
of \DIFdelbegin \DIFdel{$u'$ }\DIFdelend \DIFaddbegin \DIFadd{$\mathbf{u}'$ }\DIFaddend to $\partial M$, the third term implies
\DIFdelbegin \begin{displaymath}
\DIFdel{\left(A_{ijkl}\frac{\partial u_{k}'}{\partial x_{l}}\right)\hat{n}_{j}=h_{i}',
}\end{displaymath}%DIFAUXCMD
\DIFdelend \DIFaddbegin \begin{equation}
\DIFadd{\left(A_{ijkl}\frac{\partial u_{k}'}{\partial x_{l}}\right)\hat{n}_{j}=h_{i}',
\label{eq:27}
}\end{equation}\DIFaddend 
and finally we have the terminal conditions
\DIFdelbegin \begin{displaymath}
\DIFdel{u_{i}'(\mathbf{x},T)=\frac{\partial u_{i}'}{\partial t}(\mathbf{x},T)=0.
}\end{displaymath}%DIFAUXCMD
\DIFdelend \DIFaddbegin \begin{equation}
\DIFadd{u_{i}'(\mathbf{x},T)=\frac{\partial u_{i}'}{\partial t}(\mathbf{x},T)=0.
\label{eq:28}
}\end{equation}\DIFaddend 
It has taken some effort, but we have now shown that eq.\DIFdelbegin \DIFdel{(11}\DIFdelend \DIFaddbegin \DIFadd{~(\ref{eq:11}}\DIFaddend ) requires the Lagrange multiplier field, $\mathbf{u}'$, 
to satisfy an elastic wave equation, with this problem driven by a traction applied at the observation point \DIFdelbegin \DIFdel{and 
being }\DIFdelend \DIFaddbegin \DIFadd{that is 
}\DIFaddend equal to the difference between the predicted and observed seismograms. \DIFdelbegin \DIFdel{It has }\DIFdelend \DIFaddbegin \DIFadd{We have }\DIFaddend also shown that the second 
Lagrange multiplier field, $\mathbf{w}'$, is just the restriction of $\mathbf{u}'$ to the boundary, and 
hence no additional work is necessary for its calculation.

We can now use eq.\DIFdelbegin \DIFdel{(10}\DIFdelend \DIFaddbegin \DIFadd{~(\ref{eq:10}}\DIFaddend ) to obtain sensitivity kernels with respect to the different model parameters. For arbitrary 
variations in these parameters we have
\DIFdelbegin \begin{align*}
\DIFdel{\langle DJ(\mathbf{m})|\delta \mathbf{m}\rangle = -}&\DIFdel{\int_{0}^{T}\int_{M}\left[\delta\rho\frac{\partial^{2}u_{i}}{\partial t^{2}}-\frac{\partial}{\partial x_{j}}\left(\delta A_{ijkl}\frac{\partial u_{k}}{\partial x_{l}}-\delta\overline{S}_{ij}\right)\right]u_{i}'\dd^{3}\mathbf{x}\dd t \nonumber }\\
\DIFdel{-}&\DIFdel{\int_{0}^{T}\int_{\partial M}\left(\delta A_{ijkl}\frac{\partial u_{k}}{\partial x_{l}}-\delta\overline{S}_{ij}\right)\hat{n}_{j}u_{i}'\dd S\dd t.
}\end{align*}%DIFAUXCMD
\DIFdelend \DIFaddbegin \begin{align}
\DIFadd{\langle DJ(\mathbf{m})|\delta \mathbf{m}\rangle = -}&\DIFadd{\int_{0}^{T}\int_{M}\left[\delta\rho\frac{\partial^{2}u_{i}}{\partial t^{2}}-\frac{\partial}{\partial x_{j}}\left(\delta A_{ijkl}\frac{\partial u_{k}}{\partial x_{l}}-\delta\overline{S}_{ij}\right)\right]u_{i}'\dd^{3}\mathbf{x}\dd t \nonumber }\\
\DIFadd{-}&\DIFadd{\int_{0}^{T}\int_{\partial M}\left(\delta A_{ijkl}\frac{\partial u_{k}}{\partial x_{l}}-\delta\overline{S}_{ij}\right)\hat{n}_{j}u_{i}'\dd S\dd t.
\label{eq:29}
}\end{align}\DIFaddend 
Integrating by parts once, this expression can be simplified \DIFdelbegin \begin{displaymath}
\DIFdel{\langle DJ(\mathbf{m})|\delta \mathbf{m}\rangle=\int_{0}^{T}\int_{M}\left[\delta\rho\frac{\partial u_{i}}{\partial t}\frac{\partial u_{i}'}{\partial t}-\delta A_{ijkl}\frac{\partial u_{i}'}{\partial x_{j}}\frac{\partial u_{k}}{\partial x_{l}}+\delta\overline{S}_{ij}\frac{\partial u_{i}'}{\partial x_{j}}\right]\dd^{3}\mathbf{x}\dd t,
}\end{displaymath}%DIFAUXCMD
\DIFdelend \DIFaddbegin \DIFadd{to
}\begin{equation}
\DIFadd{\langle DJ(\mathbf{m})|\delta \mathbf{m}\rangle=\int_{0}^{T}\int_{M}\left[\delta\rho\frac{\partial u_{i}}{\partial t}\frac{\partial u_{i}'}{\partial t}-\delta A_{ijkl}\frac{\partial u_{i}'}{\partial x_{j}}\frac{\partial u_{k}}{\partial x_{l}}+\delta\overline{S}_{ij}\frac{\partial u_{i}'}{\partial x_{j}}\right]\dd^{3}\mathbf{x}\dd t,
\label{eq:30}
}\end{equation}\DIFaddend 
from which we can read off the various sensitivity kernels
\DIFdelbegin \begin{align*}
\DIFdel{K_{\rho}}&\DIFdel{=\int_{0}^{T}\frac{\partial u_{i}}{\partial t}\frac{\partial u_{i}'}{\partial t}dt, }\\
\DIFdel{K_{A_{ijkl}}}&\DIFdel{=-\int_{0}^{T}\frac{\partial u_{i}'}{\partial x_{j}}\frac{\partial u_{k}}{\partial x_{l}}dt, }\\
\DIFdel{K_{\overline{S}_{ij}}}&\DIFdel{= \frac{\partial u_{i}'}{\partial x_{j}}.
}\end{align*}%DIFAUXCMD
\DIFdelend \DIFaddbegin \begin{align}
\DIFadd{K_{\rho}(\mathbf{x})}&\DIFadd{=\int_{0}^{T}\frac{\partial u_{i}}{\partial t}(\mathbf{x},t)\frac{\partial u_{i}'}{\partial t}(\mathbf{x},t)\dd t, \label{eq:31} }\\
\DIFadd{K_{A_{ijkl}}(\mathbf{x})}&\DIFadd{=-\int_{0}^{T}\frac{\partial u_{i}'}{\partial x_{j}}(\mathbf{x},t)\frac{\partial u_{k}}{\partial x_{l}}(\mathbf{x},t)\dd t, \label{eq:32} }\\
\DIFadd{K_{\overline{S}_{ij}}(\mathbf{x},t)}&\DIFadd{= \frac{\partial u_{i}'}{\partial x_{j}}(\mathbf{x},t).
\label{eq:33}
}\end{align}\DIFaddend 
\DIFaddbegin \DIFadd{Note that the kernel for the stress glut depends on time as well as position, the stress glut itself being a function of both.
}

\DIFadd{The problem for $\mathbf{u}'$ differs from the forward problem in being subject to terminal rather than initial conditions, 
so that it must be solved backwards in time from $t=T$. This is easily dealt with. Let us define the time-reversed field
}\begin{equation}
\DIFadd{u_{i}^{\dagger}(\mathbf{x},t)=u_{i}'(\mathbf{x},T-t).
\label{eq:33a}
}\end{equation}
\DIFadd{Because eq.~(\ref{eq:26}) involves only second time derivatives, $\mathbf{u}^{\dagger}$ satisfies exactly the same wave equation, 
along with the boundary condition
}\begin{equation}
\DIFadd{\left(A_{ijkl}\frac{\partial u_{k}^{\dagger}}{\partial x_{l}}\right)\hat{n}_{j}=h_{i}^{\dagger}, \quad 
h_{i}^{\dagger}(\mathbf{x},t)=h_{i}'(\mathbf{x},T-t),
\label{eq:33b}
}\end{equation}
\DIFadd{and the terminal conditions in eq.~(\ref{eq:28}) become the initial conditions 
$u_{i}^{\dagger}(\mathbf{x},0)=\partial_{t}u_{i}^{\dagger}(\mathbf{x},0)=0$. The problem for $\mathbf{u}^{\dagger}$ is therefore an 
initial value problem of precisely the same form as the forward problem, driven at the receiver by the time-reversed residual 
$h_{i}^{\dagger}$, and it can be solved with the same numerical code. In terms of $\mathbf{u}^{\dagger}$ the sensitivity kernels read
}\begin{align}
\DIFadd{K_{\rho}(\mathbf{x})}&\DIFadd{=-\int_{0}^{T}\frac{\partial u_{i}}{\partial t}(\mathbf{x},t)\frac{\partial u_{i}^{\dagger}}{\partial t}(\mathbf{x},T-t)\dd t, \label{eq:33c} }\\
\DIFadd{K_{A_{ijkl}}(\mathbf{x})}&\DIFadd{=-\int_{0}^{T}\frac{\partial u_{i}^{\dagger}}{\partial x_{j}}(\mathbf{x},T-t)\frac{\partial u_{k}}{\partial x_{l}}(\mathbf{x},t)\dd t, \label{eq:33d} }\\
\DIFadd{K_{\overline{S}_{ij}}(\mathbf{x},t)}&\DIFadd{= \frac{\partial u_{i}^{\dagger}}{\partial x_{j}}(\mathbf{x},T-t),
\label{eq:33e}
}\end{align}
\DIFadd{where the sign change in $K_{\rho}$ comes from the time derivative of the reversed field. These expressions make the 
structure of the kernels clear: the forward field at time $t$ is paired with the field that has spread out from the receiver 
for a time $T-t$, which is just the time remaining for anything happening at $(\mathbf{x},t)$ to be felt within the observation window. 
The time reversal is nothing more than a convenient way of converting a terminal value problem into an initial value one, and 
plays no fundamental role in the theory.
}\DIFaddend 

The approach we have described is known as the \textbf{adjoint method}, and in this context 
\DIFdelbegin \DIFdel{$u_{i}'$ }\DIFdelend \DIFaddbegin \DIFadd{$\mathbf{u}^{\dagger}$ }\DIFaddend is known as the \textbf{adjoint wavefield} and \DIFdelbegin \DIFdel{$h_{i}'$ }\DIFdelend \DIFaddbegin \DIFadd{$h_{i}^{\dagger}$ }\DIFaddend the \textbf{adjoint traction} 
\DIFdelbegin \footnote{\DIFdel{In detail, it 
is conventional to define the adjoint wavefield by time-reversing $u_{i}'$ such that the terminal conditions 
obtained are transformed into initial conditions. But this step is only done for 
convenience and is not necessary.}}%DIFAUXCMD
\addtocounter{footnote}{-1}%DIFAUXCMD
\DIFdel{.}\DIFdelend \DIFaddbegin \DIFadd{or }\textbf{\DIFadd{adjoint source}}\DIFadd{.}\footnote{\DIFadd{Some authors reserve these names for the un-reversed field $\mathbf{u}'$ and 
traction $h_{i}'$. Nothing depends on this choice so long as the time arguments in the kernels are kept track of.}} 
\DIFaddend The steps needed to calculate the gradient of the misfit can be summarised as follows:
\begin{enumerate}
    \item Solve the forward equations for $\mathbf{u}$;
    \item From this solution determine \DIFdelbegin \DIFdel{J }\DIFdelend \DIFaddbegin \DIFadd{$J$ }\DIFaddend and the associated adjoint traction \DIFdelbegin \DIFdel{$h_{i}'$}\DIFdelend \DIFaddbegin \DIFadd{$h_{i}^{\dagger}$}\DIFaddend ;
    \item Solve the adjoint equations for \DIFdelbegin \DIFdel{$\mathbf{u}'$}\DIFdelend \DIFaddbegin \DIFadd{$\mathbf{u}^{\dagger}$ as an initial value problem}\DIFaddend ;
    \item Construct the sensitivity kernels by combining $\mathbf{u}$ and \DIFdelbegin \DIFdel{$\mathbf{u}'$ }\DIFdelend \DIFaddbegin \DIFadd{$\mathbf{u}^{\dagger}$ }\DIFaddend appropriately.
\end{enumerate}
The adjoint problem is identical in form to the original, and hence the cost of this method is just 
\textbf{two numerical solutions of the elastic wave equation}. In fact, within practical applications it 
is usual to perform three solutions, this being done so \DIFaddbegin \DIFadd{that }\DIFaddend the forward field $\mathbf{u}$ need not be stored at 
all times when assembling the sensitivity kernels\DIFaddbegin \DIFadd{: because the kernels pair $\mathbf{u}$ at time $t$ with 
$\mathbf{u}^{\dagger}$ at time $T-t$, the forward field is instead reconstructed backwards in time alongside the 
adjoint calculation}\DIFaddend . Even so, the saving relative to the finite-difference 
approach is by a factor of $\frac{3}{m+1}$ where $m$ is the number of model parameters. Given that $m$ in practice 
is of order \DIFdelbegin \DIFdel{$10^{4}-10^{6}$ }\DIFdelend \DIFaddbegin \DIFadd{$10^{4}$--$10^{6}$, }\DIFaddend this is an enormous improvement and justifies a few pages of algebra! 
It is the adjoint method that allows waveform inversion to be a practical tool within seismology and also 
in a range of other fields. Indeed, it is this method that underlies numerical weather prediction.

It is worth commenting that though we have obtained the adjoint equations for a specific choice of misfit, the form of $J$
 only enters into the problem through the adjoint traction $h_{i}'$ \DIFaddbegin \DIFadd{(or, equivalently, its time reverse $h_{i}^{\dagger}$)}\DIFaddend . This means that the theory can be readily applied to 
 other measures of misfit without having to go through the lengthy derivation each time.


\DIFdelbegin %DIFDELCMD < \begin{figure}[ht]
%DIFDELCMD <     \centering    
%DIFDELCMD <     \includegraphics[width=0.7\textwidth]{figures/L8F2.png}    
%DIFDELCMD <     %%%
%DIFDELCMD < \caption{%
{%DIFAUXCMD
\DIFdelFL{Images from the first global waveform tomographic model produced by Bozdag et al. (2016). On the left is shown the starting model S362ANI which was obtained using ray theoretic methods, and on the right can be seen the waveform model following 15 iterations. While the difference in the models might not seem huge, they are still significant. There is considerably more detail within the waveform model, including structures associated with a range of hot spots including that under Yellowstone. Moreover, as methods continue to improve the difference between ray theoretic and waveform tomographic models will only grow further.}}
    %DIFAUXCMD
%DIFDELCMD < \label{fig:2}
%DIFDELCMD < \end{figure}
%DIFDELCMD < %%%
\DIFdelend \DIFaddbegin \section*{\DIFadd{Banana doughnut kernels}}
\DIFadd{Nothing in the adjoint method is specific to the least-squares misfit, or indeed to misfits at all. It provides the 
derivative with respect to the model of any differentiable functional of the wavefield at the cost of one additional 
wavefield simulation, and so it is as useful for asking how a given measurement depends on Earth structure, which is 
to say for }\textbf{\DIFadd{sensitivity analysis}}\DIFadd{, as it is for the estimation of the model or the state of a system from data. 
As a first application, before returning to inversions, we use it to find the sensitivity of the most basic seismological 
measurement of all, a travel time, and compare the result with the ray-theoretic prediction of Lecture~19.
}\DIFaddend 

\DIFdelbegin \section*{\DIFdel{Banana Doughnut kernels}}
%DIFAUXCMD
\DIFdel{Let $\mathbf{u}^{obs}(t)$ }\DIFdelend \DIFaddbegin \DIFadd{Let $\mathbf{u}^{\mathrm{obs}}(t)$ }\DIFaddend again denote the displacement vector recorded by a seismograph at the surface location, 
\DIFdelbegin \DIFdel{$\mathbf{x}_{s}$}\DIFdelend \DIFaddbegin \DIFadd{$\mathbf{x}_{r}$}\DIFaddend , following an earthquake. Suppose we define
\DIFdelbegin \begin{displaymath}
\DIFdel{s^{obs}(t)=w(t)\hat{\nu}_{i}u_{i}^{obs}(t).
}\end{displaymath}%DIFAUXCMD
\DIFdelend \DIFaddbegin \begin{equation}
\DIFadd{s^{\mathrm{obs}}(t)=w(t)\hat{\nu}_{i}u_{i}^{\mathrm{obs}}(t),
\label{eq:34}
}\end{equation}\DIFaddend 
with \DIFdelbegin \DIFdel{$\hat{\boldsymbol{\nu}}$ }\DIFdelend \DIFaddbegin \DIFadd{$\hat{\bm{\nu}}$ }\DIFaddend a unit vector that selects a component of the displacement of interest, 
and $w(t)$ a windowing function that \DIFdelbegin \DIFdel{zeros }\DIFdelend \DIFaddbegin \DIFadd{zeroes }\DIFaddend everything out except for a particular seismic phase of interest. Exactly 
the same thing can be done with a synthetic seismogram
\DIFdelbegin \begin{displaymath}
\DIFdel{s(\mathbf{x}_{s},t)=w(t)\hat{\nu}_{i}u_{i}(\mathbf{x}_{s},t),
}\end{displaymath}%DIFAUXCMD
\DIFdelend \DIFaddbegin \begin{equation}
\DIFadd{s(\mathbf{x}_{r},t)=w(t)\hat{\nu}_{i}u_{i}(\mathbf{x}_{r},t),
\label{eq:35}
}\end{equation}\DIFaddend 
with $\mathbf{u}$ calculated in some reference model. Consider 
the \textbf{cross-correlation} of these two time series as defined by
\DIFdelbegin \begin{displaymath}
\DIFdel{C(\tau)=\int_{-\infty}^{\infty}s^{obs}(t-\tau)s(\mathbf{x}_{s},t)dt.
}\end{displaymath}%DIFAUXCMD
\DIFdelend \DIFaddbegin \begin{equation}
\DIFadd{C(\tau)=\int_{-\infty}^{\infty}s^{\mathrm{obs}}(t-\tau)s(\mathbf{x}_{r},t)\dd t.
\label{eq:36}
}\end{equation}\DIFaddend 
This quantity will be largest when \DIFdelbegin \DIFdel{$s^{obs}$ and s }\DIFdelend \DIFaddbegin \DIFadd{$s^{\mathrm{obs}}$ and $s$ }\DIFaddend are maximally aligned, and hence the associated value 
of $\tau$ gives a measurement of the delay time between the observed and synthetic waveforms. Indeed, this 
is typically what is done in practice to automate delay time measurements. The condition for $\overline{\tau}$ 
to maximise the cross-correlation of the waveforms is given by
\DIFdelbegin \begin{displaymath}
\DIFdel{\frac{dC}{d\tau}(\overline{\tau})=0
}\end{displaymath}%DIFAUXCMD
\DIFdelend \DIFaddbegin \begin{equation}
\DIFadd{\frac{\ddns C}{\ddns\tau}(\overline{\tau})=0,
\label{eq:37}
}\end{equation}\DIFaddend 
and this acts to define $\overline{\tau}$ implicitly as a function of the synthetic wavefield $\mathbf{u}$, and
 hence of the underlying model parameters. 

 The adjoint method can be applied to determine the functional derivatives of $\overline{\tau}$\DIFdelbegin \DIFdel{and 
 here }\DIFdelend \DIFaddbegin \DIFadd{, and 
 as remarked above }\DIFaddend we need only determine the appropriate adjoint source\DIFdelbegin \DIFdel{(details left to the interested reader) . The result is that }\DIFdelend \DIFaddbegin \DIFadd{. To do this we perturb the synthetic waveform, 
 $s\to s+\delta s$, within eq.~(\ref{eq:37}), noting that $\overline{\tau}$ changes in response. Differentiating 
 eq.~(\ref{eq:36}) gives
}\begin{equation}
\DIFadd{\frac{\ddns C}{\ddns\tau}(\tau)=-\int_{-\infty}^{\infty}\dot{s}^{\mathrm{obs}}(t-\tau)s(\mathbf{x}_{r},t)\dd t,
\label{eq:38}
}\end{equation}
\DIFadd{and so, to first order, the perturbed form of eq.~(\ref{eq:37}) reads
}\begin{equation}
\DIFadd{-\int_{-\infty}^{\infty}\dot{s}^{\mathrm{obs}}(t-\overline{\tau})\,\delta s(\mathbf{x}_{r},t)\dd t
+\delta\overline{\tau}\int_{-\infty}^{\infty}\ddot{s}^{\mathrm{obs}}(t-\overline{\tau})s(\mathbf{x}_{r},t)\dd t=0.
\label{eq:39}
}\end{equation}
\DIFadd{If the observed and synthetic waveforms differ only by the time shift, so that $s^{\mathrm{obs}}(t-\overline{\tau})=s(\mathbf{x}_{r},t)$, 
which is a reasonable approximation when the reference model is close to the truth, then an integration by parts (the 
window removing any boundary terms) gives $\int\ddot{s}s\dd t=-\int\dot{s}^{2}\dd t$, and hence
}\begin{equation}
\DIFadd{\delta\overline{\tau}=-\frac{1}{N}\int_{-\infty}^{\infty}\dot{s}(\mathbf{x}_{r},t)\,\delta s(\mathbf{x}_{r},t)\dd t,
\quad N=\int_{-\infty}^{\infty}\dot{s}(\mathbf{x}_{r},t)^{2}\dd t.
\label{eq:40}
}\end{equation}
\DIFadd{Recalling that $\delta s=w(t)\hat{\nu}_{i}\delta u_{i}(\mathbf{x}_{r},t)$, and comparing with eqs.~(\ref{eq:18}) and (\ref{eq:19}), 
we can read off the adjoint traction
}\begin{equation}
\DIFadd{h_{i}'(\mathbf{x},t)=-\frac{1}{N}w(t)\dot{s}(\mathbf{x}_{r},t)\hat{\nu}_{i}\delta(\mathbf{x}-\mathbf{x}_{r}),
\label{eq:41}
}\end{equation}
\DIFadd{which is the windowed velocity of the synthetic seismogram applied at the receiver, and the sensitivity kernels then follow 
from eqs.~(\ref{eq:31})--(\ref{eq:33}) exactly as before. Note that the observed seismogram has dropped out entirely: 
the kernel depends only on the synthetic wavefield, and hence only on the reference model. In this manner }\DIFaddend we can write
\DIFdelbegin \begin{displaymath}
\DIFdel{\delta\overline{\tau}=\int_{M}(K_{\alpha}\delta\alpha+K_{\beta}\delta\beta)\dd^{3}\mathbf{x}
}\end{displaymath}%DIFAUXCMD
\DIFdelend \DIFaddbegin \begin{equation}
\DIFadd{\delta\overline{\tau}=\int_{M}(K_{\alpha}\,\delta\alpha+K_{\beta}\,\delta\beta)\dd^{3}\mathbf{x},
\label{eq:42}
}\end{equation}\DIFaddend 
for the first-order perturbation in this cross-correlation delay time with respect to the P- and S-wave speeds, 
holding other model parameters fixed. This expression can be compared to the \DIFdelbegin \DIFdel{ray theoretic result
}\begin{displaymath}
\DIFdel{\delta T=-\int_{\text{ray}}\frac{\delta c}{c^{2}}ds,
}\end{displaymath}%DIFAUXCMD
\DIFdelend \DIFaddbegin \DIFadd{ray-theoretic result
}\begin{equation}
\DIFadd{\delta T=-\int_{\mathrm{ray}}\frac{\delta c}{c^{2}}\dd s,
\label{eq:43}
}\end{equation}\DIFaddend 
obtained using Fermat's principle, with \DIFdelbegin \DIFdel{c }\DIFdelend \DIFaddbegin \DIFadd{$c$ }\DIFaddend the appropriate phase speed. 


\DIFdelbegin %DIFDELCMD < \begin{figure}[ht]
%DIFDELCMD <     %%%
\DIFdelendFL \DIFaddbeginFL \DIFaddFL{Figs.~\ref{fig:3} and \ref{fig:4} show how such a kernel is built up in a two-dimensional numerical experiment using the 
adjoint method exactly as derived above. An explosive point source generates a P-wave that is recorded at a receiver, the 
delay time of the arrival is defined through cross-correlation, and the corresponding adjoint source is the time-reversed 
windowed velocity of the synthetic seismogram applied at the receiver. The adjoint field $\mathbf{u}^{\dagger}$ spreads out 
from the receiver as an ordinary initial value problem, and the kernel for the P-wave speed, 
}\begin{equation*}
\DIFaddFL{K_{\alpha}(\mathbf{x})=-2\rho\alpha\int_{0}^{T}\nabla\cdot\mathbf{u}(\mathbf{x},t)\,
\nabla\cdot\mathbf{u}^{\dagger}(\mathbf{x},T-t)\dd t,
}\end{equation*}
\DIFaddFL{accumulates wherever the forward field at time $t$ overlaps the adjoint field at time $T-t$. At any instant this overlap 
of the two wavefronts is confined to a pair of points, and as these points sweep from source to receiver they trace out the 
region of sensitivity, which is an ellipse with the source and receiver at its foci.}\footnote{\DIFaddFL{The kernel was checked by 
perturbing the wave speed within a small region and measuring the resulting change in the cross-correlation delay time 
directly; the kernel predicts this change to within about ten per cent for a perturbation of two per cent.}}

\begin{figure}[htbp]
    \DIFaddendFL \centering
    \DIFdelbeginFL %DIFDELCMD < \includegraphics[width=0.7\textwidth]{figures/L8F3.png}        
%DIFDELCMD <     %%%
\DIFdelendFL \DIFaddbeginFL \includegraphics[width=\textwidth]{figures/L8F3.pdf}
    \captionsetup{justification=raggedright}
    \DIFaddendFL \caption{\DIFdelbeginFL \DIFdelFL{Cross sections through the functional derivative }\DIFdelendFL \DIFaddbeginFL \DIFaddFL{Construction }\DIFaddendFL of a cross-correlation \DIFdelbeginFL \DIFdelFL{delay time for a P-wave 
    phase with respect to $\alpha$ as calculated }\DIFdelendFL \DIFaddbeginFL \DIFaddFL{delay-time kernel }\DIFaddendFL by the adjoint method \DIFaddbeginFL \DIFaddFL{in a two-dimensional homogeneous 
    body}\DIFaddendFL . \DIFdelbeginFL \DIFdelFL{Here we see that }\DIFdelendFL \DIFaddbeginFL \DIFaddFL{The star marks an explosive source and }\DIFaddendFL the \DIFdelbeginFL \DIFdelFL{delay time 
    }\DIFdelendFL \DIFaddbeginFL \DIFaddFL{triangle a receiver. The kernel is 
    \mbox{$K_{\alpha}(\mathbf{x})=-2\rho\alpha\int_{0}^{T}\nabla\cdot\mathbf{u}(\mathbf{x},t)\,\nabla\cdot\mathbf{u}^{\dagger}(\mathbf{x},T-t)\dd t$}, 
    and each row shows one instant $t$: the divergence }\DIFaddendFL of the 
    \DIFdelbeginFL \DIFdelFL{same phase can be measured at different frequencies by filtering }\DIFdelendFL \DIFaddbeginFL \DIFaddFL{forward field $\mathbf{u}(\mathbf{x},t)$; }\DIFaddendFL the \DIFdelbeginFL \DIFdelFL{data }\DIFdelendFL \DIFaddbeginFL \DIFaddFL{divergence of the adjoint field $\mathbf{u}^{\dagger}(\mathbf{x},T-t)$, 
    which has spread out from the receiver for a time $T-t$; the integrand, that is the instantaneous contribution to the 
    kernel; }\DIFaddendFL and \DIFdelbeginFL \DIFdelFL{synthetic waveforms}\DIFdelendFL \DIFaddbeginFL \DIFaddFL{the kernel accumulated over times up to $t$}\DIFaddendFL . The \DIFdelbeginFL \DIFdelFL{dominant sensitivity lies within }\DIFdelendFL \DIFaddbeginFL \DIFaddFL{integrand, and hence }\DIFaddendFL the \DIFdelbeginFL \DIFdelFL{first Fresnel zone}\DIFdelendFL \DIFaddbeginFL \DIFaddFL{kernel}\DIFaddendFL , \DIFdelbeginFL \DIFdelFL{but it curiously vanishes along }\DIFdelendFL \DIFaddbeginFL \DIFaddFL{is non-zero only where }\DIFaddendFL the \DIFdelbeginFL \DIFdelFL{ray path. 
    Figure taken from Liu \& Tromp (2008)}\DIFdelendFL \DIFaddbeginFL \DIFaddFL{two 
    fields overlap}\DIFaddendFL .}
    \label{fig:3}
\end{figure}

\DIFaddbegin \begin{figure}[htbp]
    \centering
    \includegraphics[width=\textwidth]{figures/L8F4.pdf}
    \caption{\DIFaddFL{The finished two-dimensional kernel of Fig.~\ref{fig:3} (left; the colour scale is clipped near the source 
    and receiver, where the kernel is singular) and a cross-section through it at the mid-point of the ray (right). The kernel 
    is negative throughout the first Fresnel zone, since an increase in wave speed there advances the arrival, and 
    changes sign in the surrounding zones. Unlike the three-dimensional kernels of Fig.~\ref{fig:5}, it does not vanish on 
    the ray itself, though it is somewhat smaller there than on either side.}}
    \label{fig:4}
\end{figure}

\DIFaddend The form of these sensitivity kernels deserves some comment\DIFaddbegin \DIFadd{, and Fig.~\ref{fig:5} shows examples calculated for the Earth}\DIFaddend .  First, within ray theory all sensitivity lies along 
the reference ray path. Within the adjoint theory, however, there is a broad region of sensitivity around the ray 
path. This region can be shown to be the first Fresnel zone for the wave under consideration, i.e.\DIFaddbegin \DIFadd{\ }\DIFaddend the region in 
which a scattered wave arrives in phase with the main arrival. This is not too surprising, as ray theory neglects 
diffraction effects. What is perhaps more curious is that the true kernels vanish along the reference 
ray path. When this was first observed from numerical calculations\DIFaddbegin \DIFadd{, }\DIFaddend this result seemed counter-intuitive to some 
people: ray theory says that all the sensitivity is along the ray path, but the exact theory says none of it is! 

\DIFaddbegin \begin{figure}[htbp]
    \centering
    \includegraphics[width=0.7\textwidth]{figures/L8F5.png}        
    \caption{\DIFaddFL{Cross-sections through the functional derivative of a cross-correlation delay time for a P-wave 
    phase with respect to $\alpha$ as calculated by the adjoint method. Here we see that the delay time 
    of the same phase can be measured at different frequencies by filtering the data and synthetic waveforms. 
    The dominant sensitivity lies within the first Fresnel zone, but it curiously vanishes along the ray path. 
    Figure taken from Liu \& Tromp (2008).}}
    \label{fig:5}
\end{figure}
\DIFaddend Subsequent work showed that there is no contradiction. The exact kernels vanish along the ray path because a 
scatterer on the ray path is associated with no phase delay relative to the reference wave. \DIFaddbegin \DIFadd{This argument is specific to three dimensions, 
and the two-dimensional kernel of Fig.~\ref{fig:4} has no such hole. The reason lies in the different character of wave 
propagation in even and odd numbers of dimensions, but this is a subtle matter that we will not go into. }\DIFaddend Moreover, it 
can be shown that as the dominant frequency of the wave increases, the width of the Fresnel zone shrinks, and in 
the limit as the frequency goes to infinity the two kernels agree in the sense of generalised functions,
this meaning that their action on a smooth model perturbation produces the same delay time. 

The final point to mention 
is that you might think that the kernels in Fig\DIFdelbegin \DIFdel{.3 }\DIFdelend \DIFaddbegin \DIFadd{.~\ref{fig:5} }\DIFaddend look a little like a banana cut lengthwise, while in cross 
section they are hollow and can be said to resemble a doughnut. The first person to work on this theory, 
Tony Dahlen (\DIFdelbegin \DIFdel{1942-2007}\DIFdelend \DIFaddbegin \DIFadd{1942--2007}\DIFaddend ) from Princeton, certainly thought so, and as a result he jokingly named 
them \textbf{banana doughnut kernels}, and this remains the usual term used.

\DIFaddbegin \section*{\DIFadd{Global waveform tomographic models}}
\DIFadd{An example of a global tomographic model obtained by the methods of this lecture is shown in Fig.~\ref{fig:6}.
}

\begin{figure}[htbp]
    \centering    
    \includegraphics[width=0.66\textwidth]{figures/L8F6.pdf}    
    \caption{\DIFaddFL{The global waveform tomographic model GLAD-M35 of Cui et al.\ (2024), the third generation of models produced 
    at Princeton using the adjoint method, at depths from $140\,\mathrm{km}$ to $2800\,\mathrm{km}$. The columns show the 
    perturbations in shear wave speed, bulk sound speed and compressional wave speed relative to a spherically symmetric 
    reference model. The inversion used waveforms from 2160 earthquakes recorded at 11\,575 stations and 35 iterations of 
    gradient-based optimisation in total. From Cui et al.\ (2024).}}
    \label{fig:6}
\end{figure}

\DIFadd{Such a model is the result of gradient-based 
optimisation, and as remarked earlier the question of its uncertainty is a difficult one. It can, however, be addressed 
approximately: the second derivative of the misfit at the minimum (the Hessian) plays the role of the inverse posterior 
covariance in the Bayesian picture of Lecture~19, and although it is far too large to form, its action on a vector can 
be computed with a few further forward and adjoint simulations, which is enough to estimate its largest eigenvalues and hence 
the directions in model space that the data constrain best. Fig.~\ref{fig:7} shows the resulting uncertainty estimates for the 
model of Fig.~\ref{fig:6}.
}

\begin{figure}[htbp]
    \centering    
    \includegraphics[width=0.66\textwidth]{figures/L8F7.pdf}    
    \caption{\DIFaddFL{Uncertainty in the shear wave speed of GLAD-M35 at the depths of Fig.~\ref{fig:6}: the prior standard deviation 
    assumed before inversion (left), the reduction in uncertainty implied by the data through a low-rank approximation of 
    the Hessian (centre), and the resulting posterior standard deviation (right). The uncertainty is reduced most where 
    the data coverage is best. From Cui et al.\ (2024).}}
    \label{fig:7}
\end{figure}

\DIFadd{Another way to judge a tomographic model is to compare it with models obtained by other groups using different data and methods. 
Fig.~\ref{fig:8} makes such a comparison at the base of the mantle for six models that provide both P- and S-wave speeds. 
Three of the models are the product of waveform tomography as described in this lecture, with the full non-linear inverse 
problem solved iteratively: GLAD-M35 itself, together with LOWE and its successor REVEAL, which come from a different group and 
were built with different choices of misfit, data and optimisation scheme, and with very different numbers of iterations 
(35 for GLAD-M35, 72 for LOWE and 305 for REVEAL). The other three are linearised inversions of the kind described in 
Lecture~19, in which the forward problem is linearised about a spherically symmetric reference model and solved by 
regularised least squares. GyPSuM and TX2019 are based on ray-theoretic body-wave travel times, the former combined with 
geodynamic constraints such as the gravity field and plate motions, and the latter with prior information on the positions 
of subducted slabs. SP12RTS combines travel times with surface-wave phase speeds and normal-mode observations of the kind 
discussed in Lectures~23 and~24. The models agree on the largest features, such as the two 
regions of low shear wave speed beneath Africa and the Pacific, but differ considerably in their detail and, notably, in 
the size and sign of the perturbations to the bulk sound speed, on which more will be said in Lecture~24. Such disagreements 
between models built from comparable data are a reminder that the uncertainties in tomographic models are larger than 
any single study would suggest. The same point is made by the vertical cross-sections of Fig.~\ref{fig:9}, which pass 
beneath a number of volcanic hot spots: the models differ markedly in whether they show continuous low-velocity 
conduits, or plumes, rising through the mantle beneath them.
}

\begin{figure}[htbp]
    \centering    
    \includegraphics[width=0.66\textwidth]{figures/L8F8.pdf}    
    \caption{\DIFaddFL{Six global models that provide both P- and S-wave speeds, at a depth of $2800\,\mathrm{km}$. From the top: GLAD-M35 
    (Cui et al.\ 2024), GyPSuM (Simmons et al.\ 2010), LOWE (van Herwaarden et al.\ 2020), SP12RTS (Koelemeijer et al.\ 2016), 
    TX2019 (Lu et al.\ 2019) and REVEAL (Thrastarson et al.\ 2024). The columns show the perturbations in shear wave speed, bulk 
    sound speed and compressional wave speed. From Cui et al.\ (2024).}}
    \label{fig:8}
\end{figure}

\begin{figure}[htbp]
    \centering    
    \includegraphics[width=0.85\textwidth]{figures/L8F9.pdf}    
    \caption{\DIFaddFL{Vertical cross-sections through the mantle showing the perturbation to the P-wave speed in four of the models of 
    Fig.~\ref{fig:8}: GLAD-M35, GyPSuM, TX2019 and REVEAL. The small maps show the location of each section, which in every 
    case passes beneath one or more hot spots. From Cui et al.\ (2024).}}
    \label{fig:9}
\end{figure}

\DIFaddend \section*{What you need to know and be able to do}
\DIFdelbegin %DIFDELCMD < \begin{itemize}
%DIFDELCMD <     %%%
\DIFdelend \DIFaddbegin \begin{enumerate}
    \DIFaddend \item[(i)] That waveform tomography depends on the ability to efficiently solve the elastic wave equation in realistic earth models.
    \item[(ii)] How least-squares waveform misfits can be defined, and what is meant by their gradients.
    \DIFdelbegin \DIFdel{You might also be given other waveform misfits and asked to determine the corresponding adjoint traction.
    }\DIFdelend \item[(iii)] \DIFdelbegin \DIFdel{Know how gradient based }\DIFdelend \DIFaddbegin \DIFadd{How gradient-based }\DIFaddend optimisation works, why it is useful, and what its limitations are.
    \item[(iv)] How the adjoint method can be used to calculate gradients, and what its advantages are. The derivation in this lecture is too lengthy for you to reproduce in full. But you may be asked \DIFaddbegin \DIFadd{to }\DIFaddend perform parts of the calculations with suitable guidance or to apply the method to simpler systems.
\DIFdelbegin %DIFDELCMD < \end{itemize}
%DIFDELCMD < %%%
\DIFdelend \DIFaddbegin \end{enumerate}
\DIFaddend 

\appendix
\section*{The Lagrange multiplier theorem \DIFdelbegin \DIFdel{- NON-EXAMINABLE}\DIFdelend \DIFaddbegin \DIFadd{(non-examinable)}\DIFaddend }
A sketch proof of the Lagrange multiplier theorem is provided for those interested. Let $\mathbf{u}\mapsto J(\mathbf{u})$ be a functional of the state vector, $\mathbf{u}$. Suppose that $\mathbf{u}$ is determined through solution of the following equation
\DIFdelbegin \begin{displaymath}
\DIFdel{\mathbf{a}(\mathbf{u},\mathbf{m})=\mathbf{0}.
}\end{displaymath}%DIFAUXCMD
\DIFdelend \DIFaddbegin \begin{equation}
\DIFadd{\mathbf{a}(\mathbf{u},\mathbf{m})=\mathbf{0},
\label{eq:44}
}\end{equation}\DIFaddend 
with $\mathbf{m}$ the model vector. Here \DIFdelbegin \DIFdel{$a$ }\DIFdelend \DIFaddbegin \DIFadd{$\mathbf{a}$ }\DIFaddend is some potentially non-linear vector-valued function representing abstractly the partial differential equation satisfied by $\mathbf{u}$ and in which $\mathbf{m}$ \DIFdelbegin \DIFdel{act as parameters}\DIFdelend \DIFaddbegin \DIFadd{acts as a parameter}\DIFaddend . We assume that for each $\mathbf{m}$ there is a unique solution to this equation, and hence it serves to define a mapping $\mathbf{m}\mapsto\hat{\mathbf{u}}(\mathbf{m})$ such that
\DIFdelbegin \begin{displaymath}
\DIFdel{\mathbf{a}[\hat{\mathbf{u}}(\mathbf{m}),\mathbf{m}]=\mathbf{0}.
}\end{displaymath}%DIFAUXCMD
\DIFdelend \DIFaddbegin \begin{equation}
\DIFadd{\mathbf{a}[\hat{\mathbf{u}}(\mathbf{m}),\mathbf{m}]=\mathbf{0}.
\label{eq:45}
}\end{equation}\DIFaddend 
Using this mapping, we can define a new functional
\DIFdelbegin \begin{displaymath}
\DIFdel{\hat{J}(\mathbf{m})=J}[\DIFdel{\hat{\mathbf{u}}(\mathbf{m})}]\DIFdel{.
}\end{displaymath}%DIFAUXCMD
\DIFdelend \DIFaddbegin \begin{equation}
\DIFadd{\hat{J}(\mathbf{m})=J}[\DIFadd{\hat{\mathbf{u}}(\mathbf{m})}]\DIFadd{.
\label{eq:46}
}\end{equation}\DIFaddend 
Within the main notes we do not distinguish notationally between $J(\mathbf{u})$ and $\hat{J}(\mathbf{m})$\DIFaddbegin \DIFadd{, }\DIFaddend but here it is useful to be more explicit. Our aim is to calculate the functional derivative, $D\hat{J}(\mathbf{m})$, this being the linear term within the first-order Taylor expansion
\DIFdelbegin \begin{displaymath}
\DIFdel{\hat{J}(\mathbf{m}+\delta \mathbf{m})=\hat{J}(\mathbf{m})+\langle D\hat{J}(\mathbf{m})|\delta \mathbf{m}\rangle+O(||\delta \mathbf{m}||^{2}).
}\end{displaymath}%DIFAUXCMD
\DIFdelend \DIFaddbegin \begin{equation}
\DIFadd{\hat{J}(\mathbf{m}+\delta \mathbf{m})=\hat{J}(\mathbf{m})+\langle D\hat{J}(\mathbf{m})|\delta \mathbf{m}\rangle+O(\|\delta \mathbf{m}\|^{2}).
\label{eq:47}
}\end{equation}\DIFaddend 
We will obtain this functional derivative first using a \DIFdelbegin \DIFdel{"}\DIFdelend \DIFaddbegin \DIFadd{``}\DIFaddend brute-force\DIFdelbegin \DIFdel{" }\DIFdelend \DIFaddbegin \DIFadd{'' }\DIFaddend method via the chain rule, and then show that the Lagrange multiplier theorem leads to the same result. The advantage of the latter is the relative ease of its application in practice.

Using eq.\DIFdelbegin \DIFdel{(42}\DIFdelend \DIFaddbegin \DIFadd{~(\ref{eq:46}}\DIFaddend ), we can perturb $\mathbf{m}$ and expand to \DIFdelbegin \DIFdel{first-order to obtain
}\begin{displaymath}
\DIFdel{\hat{J}(\mathbf{m}+\delta \mathbf{m})=J}[\DIFdel{\hat{\mathbf{u}}(\mathbf{m}+\delta \mathbf{m})}]\DIFdel{=
J}[\DIFdel{\hat{\mathbf{u}}(\mathbf{m})+\delta \mathbf{u}+\cdot\cdot\cdot}]\DIFdel{=\hat{J}(\mathbf{m})+
\langle DJ(\mathbf{u})|\delta \mathbf{u}\rangle+\cdot\cdot\cdot
}\end{displaymath}%DIFAUXCMD
\DIFdelend \DIFaddbegin \DIFadd{first order to obtain
}\begin{equation}
\DIFadd{\hat{J}(\mathbf{m}+\delta \mathbf{m})=J}[\DIFadd{\hat{\mathbf{u}}(\mathbf{m}+\delta \mathbf{m})}]\DIFadd{=
J}[\DIFadd{\hat{\mathbf{u}}(\mathbf{m})+\delta \mathbf{u}+\cdots}]\DIFadd{=\hat{J}(\mathbf{m})+
\langle DJ(\mathbf{u})|\delta \mathbf{u}\rangle+\cdots,
\label{eq:48}
}\end{equation}\DIFaddend 
where $\delta \mathbf{u}$ denotes the first-order perturbation to $\mathbf{u}$ due to the given change in $\mathbf{m}$. To determine the concrete form of $\delta \mathbf{u}$, we perturb eq.\DIFdelbegin \DIFdel{(41}\DIFdelend \DIFaddbegin \DIFadd{~(\ref{eq:45}}\DIFaddend ) to obtain
\DIFdelbegin \begin{displaymath}
\DIFdel{\mathbf{a}(\mathbf{u}+\delta \mathbf{u}+\cdot\cdot\cdot,\mathbf{m}+\delta \mathbf{m}+\cdot\cdot\cdot)=\mathbf{0}
}\end{displaymath}%DIFAUXCMD
\DIFdelend \DIFaddbegin \begin{equation}
\DIFadd{\mathbf{a}(\mathbf{u}+\delta \mathbf{u}+\cdots,\mathbf{m}+\delta \mathbf{m}+\cdots)=\mathbf{0},
\label{eq:49}
}\end{equation}\DIFaddend 
where it is understood that $\mathbf{u}=\hat{\mathbf{u}}(\mathbf{m})$. Expanding this out to \DIFdelbegin \DIFdel{first-order we find
}\begin{displaymath}
\DIFdel{D_{\mathbf{u}}\mathbf{a}(\mathbf{u},\mathbf{m})\delta \mathbf{u}+D_{\mathbf{m}}\mathbf{a}(\mathbf{u},\mathbf{m})\delta \mathbf{m}=\mathbf{0}
}\end{displaymath}%DIFAUXCMD
\DIFdelend \DIFaddbegin \DIFadd{first order we find
}\begin{equation}
\DIFadd{D_{\mathbf{u}}\mathbf{a}(\mathbf{u},\mathbf{m})\delta \mathbf{u}+D_{\mathbf{m}}\mathbf{a}(\mathbf{u},\mathbf{m})\delta \mathbf{m}=\mathbf{0}.
\label{eq:50}
}\end{equation}\DIFaddend 
Here we see the partial derivatives of \DIFdelbegin \DIFdel{$a$ }\DIFdelend \DIFaddbegin \DIFadd{$\mathbf{a}$ }\DIFaddend with respect to $\mathbf{u}$ and $\mathbf{m}$, and note that these are linear operators that act on the perturbations to $\mathbf{u}$ and $\mathbf{m}$ to return an appropriate vector. If this notation is confusing, think of $\mathbf{u}$ as an \DIFdelbegin \DIFdel{n-dimensional }\DIFdelend \DIFaddbegin \DIFadd{$n$-dimensional }\DIFaddend vector and $\mathbf{m}$ as an \DIFdelbegin \DIFdel{m-dimensional }\DIFdelend \DIFaddbegin \DIFadd{$m$-dimensional }\DIFaddend one. Then the values of $\mathbf{a}(\mathbf{u},\mathbf{m})$ are also \DIFdelbegin \DIFdel{n-dimensional }\DIFdelend \DIFaddbegin \DIFadd{$n$-dimensional }\DIFaddend vectors, and the linear operators $D_{\mathbf{u}}\mathbf{a}(\mathbf{u},\mathbf{m})$ and $D_{\mathbf{m}}\mathbf{a}(\mathbf{u},\mathbf{m})$ are matrices having dimensions $n\times n$ and $n\times m$, respectively. Assuming that the linear operator, $D_{\mathbf{u}}\mathbf{a}(\mathbf{u},\mathbf{m})$\DIFaddbegin \DIFadd{, }\DIFaddend is invertible, we can solve the above equation to obtain
\DIFdelbegin \begin{displaymath}
\DIFdel{\delta \mathbf{u}=-}[\DIFdel{D_{\mathbf{u}}\mathbf{a}(\mathbf{u},\mathbf{m})}]\DIFdel{^{-1}D_{\mathbf{m}}\mathbf{a}(\mathbf{u},\mathbf{m})\delta \mathbf{m}.
}\end{displaymath}%DIFAUXCMD
\DIFdelend \DIFaddbegin \begin{equation}
\DIFadd{\delta \mathbf{u}=-}[\DIFadd{D_{\mathbf{u}}\mathbf{a}(\mathbf{u},\mathbf{m})}]\DIFadd{^{-1}D_{\mathbf{m}}\mathbf{a}(\mathbf{u},\mathbf{m})\delta \mathbf{m}.
\label{eq:51}
}\end{equation}\DIFaddend 
The assumption that this matrix is invertible is not actually an assumption at all. It can be shown via the implicit function theorem that this is a necessary (and locally sufficient) condition for eq.\DIFdelbegin \DIFdel{(41}\DIFdelend \DIFaddbegin \DIFadd{~(\ref{eq:45}}\DIFaddend ) to admit unique solutions. Substituting this expression for $\delta \mathbf{u}$ into eq.\DIFdelbegin \DIFdel{(44}\DIFdelend \DIFaddbegin \DIFadd{~(\ref{eq:48}}\DIFaddend ) we find
\DIFdelbegin \begin{displaymath}
\DIFdel{\hat{J}(\mathbf{m}+\delta \mathbf{m})=\hat{J}(\mathbf{m})-\langle DJ(\mathbf{u})|
}[\DIFdel{D_{\mathbf{u}}\mathbf{a}(\mathbf{u},\mathbf{m})}]\DIFdel{^{-1}D_{\mathbf{m}}\mathbf{a}(\mathbf{u},\mathbf{m})\delta \mathbf{m}\rangle+\cdot\cdot\cdot
}\end{displaymath}%DIFAUXCMD
\DIFdelend \DIFaddbegin \begin{equation}
\DIFadd{\hat{J}(\mathbf{m}+\delta \mathbf{m})=\hat{J}(\mathbf{m})-\langle DJ(\mathbf{u})|
}[\DIFadd{D_{\mathbf{u}}\mathbf{a}(\mathbf{u},\mathbf{m})}]\DIFadd{^{-1}D_{\mathbf{m}}\mathbf{a}(\mathbf{u},\mathbf{m})\delta \mathbf{m}\rangle+\cdots.
\label{eq:52}
}\end{equation}\DIFaddend 
Recalling the definition of the adjoint of a linear operator, we can transform the inner product term to obtain
\DIFdelbegin \begin{displaymath}
\DIFdel{\hat{J}(\mathbf{m}+\delta \mathbf{m})=\hat{J}(\mathbf{m})-\langle[D_{\mathbf{m}}
\mathbf{a}(\mathbf{u},\mathbf{m})]^{\dagger}[D_{\mathbf{u}}\mathbf{a}(\mathbf{u},\mathbf{m})]^{-\dagger}DJ
(\mathbf{u})|\delta \mathbf{m}\rangle+\cdot\cdot\cdot,
}\end{displaymath}%DIFAUXCMD
\DIFdelend \DIFaddbegin \begin{equation}
\DIFadd{\hat{J}(\mathbf{m}+\delta \mathbf{m})=\hat{J}(\mathbf{m})-\langle[D_{\mathbf{m}}
\mathbf{a}(\mathbf{u},\mathbf{m})]^{\dagger}[D_{\mathbf{u}}\mathbf{a}(\mathbf{u},\mathbf{m})]^{-\dagger}DJ
(\mathbf{u})|\delta \mathbf{m}\rangle+\cdots,
\label{eq:53}
}\end{equation}\DIFaddend 
where $(\cdot)^{-\dagger}$ is a shorthand for the inverse of the adjoint. By inspection, we have arrived at the desired result
\DIFdelbegin \begin{displaymath}
\DIFdel{D\hat{J}(\mathbf{m})=-}[\DIFdel{D_{\mathbf{m}}\mathbf{a}(\mathbf{u},\mathbf{m})}]\DIFdel{^{\dagger}
}[\DIFdel{D_{\mathbf{u}}\mathbf{a}(\mathbf{u},\mathbf{m})}]\DIFdel{^{-\dagger}DJ(\mathbf{u}).
}\end{displaymath}%DIFAUXCMD
\DIFdelend \DIFaddbegin \begin{equation}
\DIFadd{D\hat{J}(\mathbf{m})=-}[\DIFadd{D_{\mathbf{m}}\mathbf{a}(\mathbf{u},\mathbf{m})}]\DIFadd{^{\dagger}
}[\DIFadd{D_{\mathbf{u}}\mathbf{a}(\mathbf{u},\mathbf{m})}]\DIFadd{^{-\dagger}DJ(\mathbf{u}).
\label{eq:54}
}\end{equation}\DIFaddend 
We now apply the Lagrange multiplier theorem and show that it leads to an identical result. The starting point is defining the appropriate Lagrangian
\DIFdelbegin \begin{displaymath}
\DIFdel{L(\mathbf{u},\mathbf{u}',\mathbf{m})=J(\mathbf{u})+\langle \mathbf{a}(\mathbf{u},\mathbf{m})|\mathbf{u}'\rangle
}\end{displaymath}%DIFAUXCMD
\DIFdelend \DIFaddbegin \begin{equation}
\DIFadd{L(\mathbf{u},\mathbf{u}',\mathbf{m})=J(\mathbf{u})+\langle \mathbf{a}(\mathbf{u},\mathbf{m})|\mathbf{u}'\rangle,
\label{eq:55}
}\end{equation}\DIFaddend 
with $\mathbf{u}'$ the Lagrange multiplier field used to impose eq.\DIFdelbegin \DIFdel{(41}\DIFdelend \DIFaddbegin \DIFadd{~(\ref{eq:45}}\DIFaddend ). Our claim is that
\DIFdelbegin \begin{displaymath}
\DIFdel{D\hat{J}(\mathbf{m})=D_{\mathbf{m}}L(\mathbf{u},\mathbf{u}',\mathbf{m}),
}\end{displaymath}%DIFAUXCMD
\DIFdelend \DIFaddbegin \begin{equation}
\DIFadd{D\hat{J}(\mathbf{m})=D_{\mathbf{m}}L(\mathbf{u},\mathbf{u}',\mathbf{m}),
\label{eq:56}
}\end{equation}\DIFaddend 
so long as the following conditions hold
\DIFdelbegin \begin{displaymath}
\DIFdel{D_{\mathbf{u}}L(\mathbf{u},\mathbf{u}',\mathbf{m})=0, \quad D_{\mathbf{u}'}L(\mathbf{u},\mathbf{u}',\mathbf{m})=0.
}\end{displaymath}%DIFAUXCMD
\DIFdelend \DIFaddbegin \begin{equation}
\DIFadd{D_{\mathbf{u}}L(\mathbf{u},\mathbf{u}',\mathbf{m})=0, \quad D_{\mathbf{u}'}L(\mathbf{u},\mathbf{u}',\mathbf{m})=0.
\label{eq:57}
}\end{equation}\DIFaddend 
As in the main lecture, we readily see that the second condition implies eq.\DIFdelbegin \DIFdel{(41}\DIFdelend \DIFaddbegin \DIFadd{~(\ref{eq:45}}\DIFaddend ), and hence we have
 $\mathbf{u}=\hat{\mathbf{u}}(\mathbf{m})$. To examine the first condition, we note that
\DIFdelbegin \begin{displaymath}
\DIFdel{\langle D_{\mathbf{u}}L(\mathbf{u},\mathbf{u}',\mathbf{m})|\delta \mathbf{u}\rangle=\langle 
    DJ(\mathbf{u})|\delta \mathbf{u}\rangle+\langle D_{\mathbf{u}}\mathbf{a}(\mathbf{u},\mathbf{m})\delta \mathbf{u}|\mathbf{u}'\rangle=0
}\end{displaymath}%DIFAUXCMD
\DIFdelend \DIFaddbegin \begin{equation}
\DIFadd{\langle D_{\mathbf{u}}L(\mathbf{u},\mathbf{u}',\mathbf{m})|\delta \mathbf{u}\rangle=\langle 
    DJ(\mathbf{u})|\delta \mathbf{u}\rangle+\langle D_{\mathbf{u}}\mathbf{a}(\mathbf{u},\mathbf{m})\delta \mathbf{u}|\mathbf{u}'\rangle=0,
\label{eq:58}
}\end{equation}\DIFaddend 
for any $\delta \mathbf{u}$. Again using the adjoint of a linear operator, this is equivalent to
\DIFdelbegin \begin{displaymath}
\DIFdel{\langle DJ(\mathbf{u})+}[\DIFdel{D_{\mathbf{u}}\mathbf{a}(\mathbf{u},\mathbf{m})}]\DIFdel{^{\dagger}\mathbf{u}'|\delta \mathbf{u}\rangle=0
}\end{displaymath}%DIFAUXCMD
\DIFdelend \DIFaddbegin \begin{equation}
\DIFadd{\langle DJ(\mathbf{u})+}[\DIFadd{D_{\mathbf{u}}\mathbf{a}(\mathbf{u},\mathbf{m})}]\DIFadd{^{\dagger}\mathbf{u}'|\delta \mathbf{u}\rangle=0,
\label{eq:59}
}\end{equation}\DIFaddend 
and as $\delta \mathbf{u}$ is arbitrary, we see that the first condition in eq.\DIFdelbegin \DIFdel{(53) implies
}\begin{displaymath}
\DIFdel{DJ(\mathbf{u})+}[\DIFdel{D_{\mathbf{u}}\mathbf{a}(\mathbf{u},\mathbf{m})}]\DIFdel{^{\dagger}\mathbf{u}'=\mathbf{0}.
}\end{displaymath}%DIFAUXCMD
\DIFdelend \DIFaddbegin \DIFadd{~(\ref{eq:57}) implies
}\begin{equation}
\DIFadd{DJ(\mathbf{u})+}[\DIFadd{D_{\mathbf{u}}\mathbf{a}(\mathbf{u},\mathbf{m})}]\DIFadd{^{\dagger}\mathbf{u}'=\mathbf{0}.
\label{eq:60}
}\end{equation}\DIFaddend 
This is a linear equation for the Lagrange multiplier $\mathbf{u}'$, and it has solution
\DIFdelbegin \begin{displaymath}
\DIFdel{\mathbf{u}'=-}[\DIFdel{D_{\mathbf{u}}\mathbf{a}(\mathbf{u},\mathbf{m})}]\DIFdel{^{-\dagger}DJ(\mathbf{u}).
}\end{displaymath}%DIFAUXCMD
\DIFdelend \DIFaddbegin \begin{equation}
\DIFadd{\mathbf{u}'=-}[\DIFadd{D_{\mathbf{u}}\mathbf{a}(\mathbf{u},\mathbf{m})}]\DIFadd{^{-\dagger}DJ(\mathbf{u}).
\label{eq:61}
}\end{equation}\DIFaddend 
This result can be substituted into eq.\DIFdelbegin \DIFdel{(52}\DIFdelend \DIFaddbegin \DIFadd{~(\ref{eq:56}}\DIFaddend ), but first we note that
\DIFdelbegin \begin{displaymath}
\DIFdel{\langle D_{\mathbf{m}}L(\mathbf{u},\mathbf{u}',\mathbf{m})|\delta \mathbf{m}\rangle=\langle D_{\mathbf{m}}\mathbf{a}(\mathbf{u},\mathbf{m})\delta \mathbf{m}|\mathbf{u}'\rangle=\langle\delta \mathbf{m}|}[\DIFdel{D_{\mathbf{m}}\mathbf{a}(\mathbf{u},\mathbf{m})}]\DIFdel{^{\dagger}\mathbf{u}'\rangle
}\end{displaymath}%DIFAUXCMD
\DIFdelend \DIFaddbegin \begin{equation}
\DIFadd{\langle D_{\mathbf{m}}L(\mathbf{u},\mathbf{u}',\mathbf{m})|\delta \mathbf{m}\rangle=\langle D_{\mathbf{m}}\mathbf{a}(\mathbf{u},\mathbf{m})\delta \mathbf{m}|\mathbf{u}'\rangle=\langle\delta \mathbf{m}|}[\DIFadd{D_{\mathbf{m}}\mathbf{a}(\mathbf{u},\mathbf{m})}]\DIFadd{^{\dagger}\mathbf{u}'\rangle,
\label{eq:62}
}\end{equation}\DIFaddend 
for any $\delta \mathbf{m}$, and hence
\DIFdelbegin \begin{displaymath}
\DIFdel{D_{\mathbf{m}}L(\mathbf{u},\mathbf{u}',\mathbf{m})=}[\DIFdel{D_{\mathbf{m}}\mathbf{a}(\mathbf{u},\mathbf{m})}]\DIFdel{^{\dagger}\mathbf{u}'=-}[\DIFdel{D_{\mathbf{m}}\mathbf{a}(\mathbf{u},\mathbf{m})}]\DIFdel{^{\dagger}[D_{\mathbf{u}}\mathbf{a}(\mathbf{u},\mathbf{m})]^{-\dagger}DJ(\mathbf{u}).
}\end{displaymath}%DIFAUXCMD
\DIFdelend \DIFaddbegin \begin{equation}
\DIFadd{D_{\mathbf{m}}L(\mathbf{u},\mathbf{u}',\mathbf{m})=}[\DIFadd{D_{\mathbf{m}}\mathbf{a}(\mathbf{u},\mathbf{m})}]\DIFadd{^{\dagger}\mathbf{u}'=-}[\DIFadd{D_{\mathbf{m}}\mathbf{a}(\mathbf{u},\mathbf{m})}]\DIFadd{^{\dagger}[D_{\mathbf{u}}\mathbf{a}(\mathbf{u},\mathbf{m})]^{-\dagger}DJ(\mathbf{u}).
\label{eq:63}
}\end{equation}\DIFaddend 
The final result is precisely what was obtained using the direct method, and hence we have established the validity of the Lagrange multiplier theorem. It is worth noting that eq.\DIFdelbegin \DIFdel{(56}\DIFdelend \DIFaddbegin \DIFadd{~(\ref{eq:60}}\DIFaddend ) for the Lagrange multiplier, $\mathbf{u}'$, involves the adjoint of the linearised forward operator. It is this fact that gives the adjoint method its name.

\end{document}
